\documentclass[11pt,a4paper]{article}

\usepackage[T1]{fontenc}
\usepackage[utf8]{inputenc}
\usepackage{amsmath,amssymb,amsthm,mathrsfs}
\usepackage{graphicx,subcaption,booktabs}
\usepackage[dvipsnames]{xcolor}
\usepackage[margin=2.6cm]{geometry}
\usepackage[colorlinks,linkcolor=MidnightBlue,citecolor=MidnightBlue,
            urlcolor=MidnightBlue]{hyperref}
\usepackage[protrusion=true,expansion=false]{microtype}

\renewcommand{\proofname}{Demonstra\c{c}\~ao}
\theoremstyle{plain}
\newtheorem{teorema}{Teorema}[section]
\newtheorem{proposicao}[teorema]{Proposição}
\newtheorem{lema}[teorema]{Lema}
\newtheorem{corolario}[teorema]{Corolário}
\theoremstyle{definition}
\newtheorem{definicao}[teorema]{Definição}
\newtheorem{exemplo}[teorema]{Exemplo}
\theoremstyle{remark}
\newtheorem{observacao}[teorema]{Observação}

\newcommand{\C}{\mathbb{C}}
\newcommand{\R}{\mathbb{R}}
\newcommand{\Z}{\mathbb{Z}}
\newcommand{\Q}{\mathbb{Q}}
\newcommand{\Zi}{\Z[i]}
\newcommand{\Qq}{\mathbf{Q}}
\newcommand{\FF}{\mathscr{F}}
\newcommand{\wpL}{\wp}
\DeclareMathOperator{\Aut}{Aut}
\DeclareMathOperator{\End}{End}
\DeclareMathOperator{\covol}{covol}
\DeclareMathOperator{\Id}{Id}

\title{\bfseries Funções complexas flexionáveis\\[2pt]
\large Tetraflexágonos quadrados, grupos cristalográficos planos\\
e funções elípticas com multiplicação complexa}
\author{Gabriel Ribeiro\\[-2pt]
\small Departamento de Matem\'atica, UFMG}
\date{\small Agosto de 2026}

\begin{document}
\maketitle
\vspace{-2.2\baselineskip}

\begin{abstract}
\noindent
Um \emph{flex} de uma face ladrilhada por polígonos congruentes é uma bijeção
que é isometria direta em cada ladrilho e os permuta --- o rearranjo que um
flexágono impõe à figura de uma face quando ela reaparece. Dada $f$ meromorfa,
perguntamos quando a imagem flexionada $\Phi_*f=f\circ\Phi^{-1}$ ainda é uma
função meromorfa. Mostramos que isso ocorre se e somente se $f$ é invariante
pelo grupo $\Gamma_\Phi=\langle A_k^{-1}A_j\rangle$ gerado pelas discrepâncias
entre as peças de $\Phi$, e que $\Gamma_\Phi$ é sempre um subgrupo discreto de
$C_n\ltimes\frac1d\mathcal O_n$, com $\mathcal O_4=\Zi$ e
$\mathcal O_3=\mathcal O_6=\Z[\zeta_3]$. O problema cai assim inteiramente
dentro da teoria das funções automorfas planas e admite classificação
completa: nove classes, cada uma com um \emph{Hauptmodul} explícito. Quando o
reticulado de $\Gamma_\Phi$ tem posto $2$ não há solução inteira não constante
e a curva $\C/\Lambda$ tem multiplicação complexa --- \emph{lemniscática}
($j=1728$) no caso quadrado, \emph{equianarmônica} ($j=0$) no
triangular. Enumeramos exaustivamente os dois grupos de flexes ($6144$ e
$524\,880$ elementos) e resolvemos os quatro flexágonos clássicos:
o tri-tetraflexágono dá $h(e^{i\pi z})$ (com soluções inteiras), o
hexa-tetraflexágono dá $\C(\wp,\wp')$ do reticulado quadrado $2\Zi$, e os
dois hexaflexágonos --- confinados pelo \emph{teorema das retas de dobra} às
rotações em torno do centro do hexágono --- dão ambos o grupo \emph{finito}
$C_3$, com soluções $f(z)=h\big((z-O)^3\big)$, entre elas inteiras. Um modelo
de flexão que acompanha a ordem das camadas decide a última questão em aberto,
reproduz exatamente o diagrama de flexão do brinquedo, e o grafo desse modelo
é decomposto por inteiro: $20\,405$ componentes, uma delas --- o flexágono ---
com $19\,200$ estados e todas as outras com $200$ ou menos. Acompanha o texto um
programa que gera e certifica a família completa de soluções de qualquer
flexágono, e catálogos prontos para impressão.
\end{abstract}

\medskip
\noindent\textbf{Palavras-chave:} flexágonos, retratos de fase, grupos
cristalográficos planos, funções elípticas, multiplicação complexa, curva
lemniscática.

\section{Introdução}

Um flexágono é um objeto de papel que, dobrado adequadamente, exibe faces
escondidas. Nos \emph{tetraflexágonos quadrados} cada face é um quadrado
subdividido em quatro quadradinhos, e o que se observa ao flexionar é que os
quatro quadradinhos de uma face reaparecem \emph{permutados e girados}. Se
pintamos uma face com o retrato de fase de uma função complexa $f$, a flexão
devolve um novo desenho: as quatro peças de $f$ recolocadas de outro modo.
Quase sempre o resultado é uma figura quebrada --- as cores não casam nos cortes.
A pergunta que motivou este trabalho é:

\begin{quote}
\emph{Para quais $f$ o desenho flexionado é, ainda, o retrato de fase de uma
função complexa?}
\end{quote}

A resposta tem três camadas. A primeira é uma tradução: o rearranjo é uma
aplicação afim por partes $\Phi$, e a condição é que as quatro determinações
$f\circ A_j^{-1}$ colem, o que pelo princípio da identidade equivale a uma
\emph{equação funcional} --- $f$ deve ser invariante por um grupo $\Gamma_\Phi$
de isometrias diretas do plano (Teorema~\ref{thm:principal}). A segunda é
estrutural: $\Gamma_\Phi$ é sempre discreto e contido em
$\mu_4\ltimes\frac12\Zi$, isto é, é um subgrupo de um grupo cristalográfico do
tipo $p4$; assim o problema cai inteiramente dentro da teoria das funções
automorfas planas e admite classificação completa
(Teorema~\ref{thm:classificacao}). A terceira é aritmética e um pouco
surpreendente: quando $\Gamma_\Phi$ contém duas translações independentes ---
o que acontece em $92\%$ dos casos --- não há solução inteira não constante, as
soluções são elípticas, e o reticulado, por estar contido em $\frac12\Zi$, é
comensurável com $\Zi$. \emph{O flexágono força multiplicação complexa}
(Corolário~\ref{cor:cm}).

Nos flexágonos concretos do manual \cite{manual} tudo se calcula: o
tri-tetraflexágono dá $h(e^{i\pi z})$ (Teorema~\ref{thm:tri}) e o hexa dá as
funções elípticas do reticulado quadrado $2\Zi$ (Teorema~\ref{thm:hexa}) --- a
função de Weierstrass do caso lemniscático, $g_3=0$, $j=1728$. A diferença
entre os dois é uma face que reaparece ao longo de dois eixos em vez de um.

\medskip
As Seções \ref{sec:flex}--\ref{sec:classificacao} são independentes de
flexágonos concretos. A Seção~\ref{sec:flexagonos} determina, por enumeração
exata das dobraduras planas, quais flexes o hexa-tetraflexágono de \cite{manual}
realmente realiza; a Seção~\ref{sec:principais} contém o teorema específico. A
Seção~\ref{sec:naoanalitico} mostra que toda a rigidez é um fenômeno
estritamente holomorfo, a Seção~\ref{sec:gerador} descreve o programa que gera a
família completa de soluções e a Seção~\ref{sec:computacional} a verificação
computacional.
\section{Flexes de quadrantes}\label{sec:flex}

Escrevemos $Q=[-1,1]^2\subset\C$ e denotamos por
\[
Q_1=[0,1]^2,\quad Q_2=[-1,0]\times[0,1],\quad
Q_3=[-1,0]^2,\quad Q_4=[0,1]\times[-1,0]
\]
os quadrantes, numerados no sentido cartesiano (a mesma convenção de
\cite[\S3.4]{manual}), com centros
\[
c_1=\tfrac{1+i}{2},\qquad c_2=\tfrac{-1+i}{2},\qquad
c_3=\tfrac{-1-i}{2},\qquad c_4=\tfrac{1-i}{2}.
\]
Seja $\mu_4=\{1,i,-1,-i\}$.

\begin{definicao}\label{def:flex}
Um \emph{flex de quadrantes} é uma aplicação $\Phi:Q\to Q$, bijetora fora dos
eixos coordenados, para a qual existem $\sigma\in S_4$ e
$\varepsilon_1,\dots,\varepsilon_4\in\mu_4$ tais que
\begin{equation}\label{eq:flex}
\Phi|_{\mathring Q_j}=A_j,\qquad
A_j(z):=\varepsilon_j\,(z-c_j)+c_{\sigma(j)}
      =\varepsilon_j z+\beta_j,\quad
\beta_j:=c_{\sigma(j)}-\varepsilon_j c_j .
\end{equation}
Denotamos $\Phi=(\sigma;\varepsilon)$ e escrevemos $\FF$ para o conjunto de
todos os flexes de quadrantes.
\end{definicao}

Cada $A_j$ é uma isometria direta que leva $Q_j$ sobre $Q_{\sigma(j)}$; a
restrição $\varepsilon_j\in\mu_4$ é forçada pelo fato de $A_j$ levar um quadrado
do reticulado em outro. A exclusão de isometrias \emph{inversas} corresponde ao
fato físico de que, numa dada face do flexágono, todos os quadradinhos são
vistos pelo mesmo lado do papel (veja a Observação~\ref{obs:reflexao}).

\begin{proposicao}
$\FF$ é um grupo para a composição, isomorfo ao produto entrelaçado
$\mu_4\wr S_4=\mu_4^{\,4}\rtimes S_4$; em particular
$|\FF|=4^4\cdot 4!=6144$.
\end{proposicao}

\begin{proof}
A aplicação $\Phi\mapsto(\varepsilon;\sigma)$ é bijetiva por
\eqref{eq:flex}, e $(\sigma;\varepsilon)\circ(\sigma';\varepsilon')
=(\sigma\sigma';\,(\varepsilon_{\sigma'(j)}\varepsilon'_j)_j)$, que é
exatamente a lei do produto entrelaçado.
\end{proof}

\begin{figure}[t]
\centering
\includegraphics[width=\linewidth]{fig/flexao_tetra.png}
\caption{A flexão de um tetraflexágono, o análogo quadrado da pinça. (a) a face
e os dois cortes; (b) dobra-se por uma mediana e os quatro quadrantes viram
duas pilhas; (c) reabre-se pela \emph{mesma} mediana, mas separando a pilha
noutro ponto: sobem folhas que estavam escondidas e aparece outra face --- são
literalmente as operações \texttt{dobrar} e \texttt{abrir} do modelo com
camadas da Seção~\ref{sec:camadas}; (d) o mapa de retorno diagonal da face $1$,
$\sigma=(13)(24)$ com $\varepsilon\equiv1$: os quadrantes opostos trocam de
lugar sem girar; (e) as discrepâncias $A_k^{-1}A_j$ desse mapa são as
translações por $2$ e por $2i$, e geram
$\Lambda_\Phi=2\Zi$. Comparar com a Figura~\ref{fig:pinca}, onde o mesmo
raciocínio no hexágono só produz rotações.}\label{fig:flexaotetra}
\end{figure}

\begin{definicao}
Dada $f$ meromorfa, o \emph{flexionado} de $f$ por $\Phi$ é
$\Phi_*f:=f\circ\Phi^{-1}$, definido em $Q$ a menos dos eixos: em
$\mathring Q_{\sigma(j)}$ vale $\Phi_*f=f\circ A_j^{-1}$.
Dizemos que $f$ é \emph{$\Phi$-flexionável} se existe uma vizinhança aberta e
conexa $U\supset Q$ e uma função $G$ meromorfa em $U$ com
$G|_{\mathring Q_{\sigma(j)}}=f\circ A_j^{-1}$ para $j=1,2,3,4$.
\end{definicao}

Note que $\Phi_*f$ é automaticamente holomorfa (meromorfa) no \emph{interior de
cada quadrante}, pois ali é composição de $f$ com uma aplicação afim. Todo o
conteúdo do problema está na colagem ao longo dos dois cortes.

\section{O teorema de caracterização}\label{sec:teorema}

\begin{definicao}\label{def:gamma}
O \emph{grupo de discrepância} de $\Phi=(\sigma;\varepsilon)$ é
\[
\Gamma_\Phi:=\bigl\langle\, A_k^{-1}\!\circ A_j \ :\ 1\le j,k\le 4 \,\bigr\rangle
\ \le\ \Aut_{\mathrm{af}}(\C),
\qquad
A_k^{-1}A_j(z)=\frac{\varepsilon_j}{\varepsilon_k}\,z
 +\frac{\beta_j-\beta_k}{\varepsilon_k}.
\]
\end{definicao}

\begin{teorema}[Caracterização]\label{thm:principal}
Seja $f$ meromorfa em $\C$ e $\Phi=(\sigma;\varepsilon)\in\FF$.
São equivalentes:
\begin{enumerate}
\item[(i)] $f$ é $\Phi$-flexionável;
\item[(ii)] as quatro determinações coincidem:
      $f\circ A_1^{-1}=f\circ A_2^{-1}=f\circ A_3^{-1}=f\circ A_4^{-1}$ em $\C$;
\item[(iii)] $f\circ S=f$ para todo $S\in\Gamma_\Phi$.
\end{enumerate}
Nesse caso $\Phi_*f=f\circ A_1^{-1}$ é meromorfa em $\C$, e o seu grupo de
invariância é $A_1\Gamma_\Phi A_1^{-1}$.
\end{teorema}

\begin{proof}
(i)$\Rightarrow$(ii). Seja $G$ meromorfa em $U\supset Q$ conexo com
$G=f\circ A_j^{-1}$ em $\mathring Q_{\sigma(j)}$. Como $A_j$ é afim inversível,
$f\circ A_j^{-1}$ é meromorfa em todo $\C$; ela coincide com $G$ no aberto não
vazio $\mathring Q_{\sigma(j)}\subset U$, e $U$ é conexo, logo pelo princípio da
identidade $G=f\circ A_j^{-1}$ em $U$. Isso vale para os quatro índices, donde
$f\circ A_j^{-1}=f\circ A_k^{-1}$ em $U$ e, de novo pelo princípio da
identidade (agora em $\C$), em todo o plano.

(ii)$\Rightarrow$(iii). De $f\circ A_j^{-1}=f\circ A_k^{-1}$ segue, compondo à
direita com $A_j$, que $f=f\circ A_k^{-1}A_j$. Os $A_k^{-1}A_j$ geram
$\Gamma_\Phi$ e o conjunto $\{S: f\circ S=f\}$ é um subgrupo.

(iii)$\Rightarrow$(i). Se $f$ é $\Gamma_\Phi$-invariante então
$f\circ A_k^{-1}A_j=f$, isto é, $f\circ A_j^{-1}=f\circ A_k^{-1}$; tome
$G:=f\circ A_1^{-1}$, meromorfa em $\C$. A última afirmação é imediata.
\end{proof}

\begin{observacao}[versão local, via Morera]\label{obs:local}
Se $f$ é apenas holomorfa numa vizinhança de $Q$ e exigimos somente que
$\Phi_*f$ se estenda \emph{continuamente} aos cortes, o teorema de
Morera--Painlevé garante que a extensão é holomorfa, e o argumento acima se
aplica localmente: como os quatro quadrantes formam um ciclo de comprimento $4$
no grafo de adjacência, basta a colagem entre pares adjacentes para concluir que
todas as determinações coincidem. Obtém-se a mesma equação funcional, agora
válida na órbita de $U$ sob $\Gamma_\Phi$; se essa órbita for todo o plano, $f$
se estende a uma função $\Gamma_\Phi$-invariante em $\C$.
\end{observacao}

\begin{observacao}[por que só isometrias diretas]\label{obs:reflexao}
Se admitíssemos que algumas peças de $\Phi$ fossem isometrias \emph{inversas}
(o que corresponderia a ver alguns quadradinhos pelo verso), a peça
correspondente de $\Phi_*f$ seria antiholomorfa. Uma função simultaneamente
holomorfa e antiholomorfa num aberto é constante; logo, num flex misto,
$\Phi_*f$ meromorfa obriga $f$ a ser constante. A hipótese de
Definição~\ref{def:flex} não é, portanto, uma restrição arbitrária: é a única
situação com conteúdo.
\end{observacao}

O próximo lema mostra que $\Gamma_\Phi$ não vê rotações globais --- em
particular, ambiguidades de convenção sobre ``como o flexágono é segurado''
são irrelevantes.

\begin{lema}\label{lem:rho}
Se $\rho$ é uma isometria direta de $\C$ com $\rho(Q)=Q$ (isto é,
$\rho(z)=i^u z$, $u\in\Z/4$), então $\rho\circ\Phi\in\FF$ e
$\Gamma_{\rho\circ\Phi}=\Gamma_\Phi$. Mais geralmente,
$\Gamma_\Phi$ depende apenas da família $(A_j)_j$ a menos de um fator comum
à esquerda.
\end{lema}

\begin{proof}
As peças de $\rho\circ\Phi$ são $\rho A_j$, e
$(\rho A_k)^{-1}(\rho A_j)=A_k^{-1}\rho^{-1}\rho A_j=A_k^{-1}A_j$.
\end{proof}

\begin{proposicao}\label{prop:discreto}
Para todo $\Phi\in\FF$ tem-se
$\Gamma_\Phi\le \mu_4\ltimes\tfrac12\Zi$. Em particular $\Gamma_\Phi$ é um
subgrupo discreto do grupo das isometrias diretas de $\C$, isto é, um grupo
cíclico finito de rotações, um grupo de friso direto, ou um grupo
cristalográfico plano dos tipos $p1$, $p2$ ou $p4$; e o seu reticulado de
translações $\Lambda_\Phi$ satisfaz $\Lambda_\Phi\subset\frac12\Zi$.
\end{proposicao}

\begin{proof}
Como $2c_j\in\Zi$ e $\varepsilon_j\in\mu_4$, temos $2\beta_j\in\Zi$; logo
$A_j$ tem a forma $z\mapsto \varepsilon_j z+t_j/2$ com $t_j\in\Zi$. O conjunto
$\{z\mapsto \varepsilon z+t/2:\varepsilon\in\mu_4,\ t\in\Zi\}$ é um grupo
(pois $\mu_4\Zi=\Zi$), obviamente discreto, e contém todos os $A_k^{-1}A_j$.
Um subgrupo discreto do grupo das isometrias diretas do plano tem parte de
translações de posto $0$, $1$ ou $2$ e grupo de pontos cíclico; aqui o grupo de
pontos é um subgrupo de $\mu_4$, o que exclui $p3$ e $p6$.
\end{proof}

\section{Classificação completa}\label{sec:classificacao}

Escrevemos $\Lambda=\Lambda_\Phi$ para o reticulado de translações de
$\Gamma_\Phi$ e $P=P_\Phi\le\mu_4$ para o grupo de pontos, de modo que
$1\to\Lambda\to\Gamma_\Phi\to P\to 1$. Quando $P\ne\{1\}$, fixamos um centro de
rotação $p$ de ordem máxima.

\begin{teorema}[Classificação]\label{thm:classificacao}
Seja $\Phi\in\FF$ e $f$ meromorfa não constante. Então $f$ é $\Phi$-flexionável
se e somente se $f$ pertence ao corpo indicado na tabela abaixo, onde
$n=|P|$, $\omega$ gera $\Lambda$ no caso de posto $1$, e $\wp=\wp(\,\cdot\,;\Lambda)$.
\begin{center}
\renewcommand{\arraystretch}{1.25}
\begin{tabular}{ccll}
\toprule
$\operatorname{posto}\Lambda$ & $n$ & corpo das soluções & inteiras não constantes?\\
\midrule
$0$ & $1$ & todas as funções meromorfas & sim\\
$0$ & $2$ & $\{h((z-p)^2)\}$ & sim\\
$0$ & $4$ & $\{h((z-p)^4)\}$ & sim\\
$1$ & $1$ & $\{h(e^{2\pi i (z-p)/\omega})\}$ & sim\\
$1$ & $2$ & $\{H(\cos(2\pi (z-p)/\omega))\}$ & sim\\
$2$ & $1$ & $\C(\wp,\wp')$ & \textbf{não}\\
$2$ & $2$ & $\C\bigl(\wp(z-p)\bigr)$ & \textbf{não}\\
$2$ & $4$ & $\C\bigl(\wp(z-p)^2\bigr)$, $\Lambda$ quadrado & \textbf{não}\\
\bottomrule
\end{tabular}
\end{center}
(Nos casos de posto $\le 1$, $h$ e $H$ percorrem todas as funções meromorfas do
argumento indicado.)
\end{teorema}

\begin{proof}
Pelo Teorema~\ref{thm:principal} basta descrever o corpo
$\mathscr{M}(\C)^{\Gamma_\Phi}$ das funções meromorfas $\Gamma_\Phi$-invariantes,
isto é, das funções meromorfas na órbifold $\C/\Gamma_\Phi$.

\emph{Posto $0$.} $\Gamma_\Phi=\langle z\mapsto \zeta_n(z-p)+p\rangle\cong C_n$ e
$w\mapsto w^n$ realiza $\C\to\C/C_n$; logo as invariantes são $h((z-p)^n)$.

\emph{Posto $1$, $n=1$.} $\Gamma_\Phi=\omega\Z$ e $z\mapsto e^{2\pi i(z-p)/\omega}$
identifica $\C/\Gamma_\Phi\cong\C^\ast$.

\emph{Posto $1$, $n=2$.} O grupo de pontos age no reticulado de posto $1$, logo
$n\le 2$, e $n=2$ significa $\Gamma_\Phi=\langle t_\omega, r_p\rangle$ com
$r_p(z)=2p-z$ e $r_pt_\omega r_p^{-1}=t_{-\omega}$. Pondo
$u=e^{2\pi i (z-p)/\omega}$, o grupo torna-se $u\mapsto u^{\pm1}$ em $\C^\ast$,
cujo quociente é $\C$ via $u+u^{-1}=2\cos(2\pi(z-p)/\omega)$.

\emph{Posto $2$, $n=1$.} Definição de função elíptica; o corpo é $\C(\wp,\wp')$.

\emph{Posto $2$, $n=2$.} As invariantes são as elípticas pares em torno de $p$,
que formam $\C(\wp(z-p))$.

\emph{Posto $2$, $n=4$.} $n=4$ obriga $i\Lambda=\Lambda$, isto é, $\Lambda$
quadrado; então $g_3(\Lambda)=0$ e, com $w=z-p$,
\begin{equation}\label{eq:CM}
\wp(iw)=-\wp(w),\qquad \wp'(iw)=i\,\wp'(w).
\end{equation}
Logo $\wp^2$ é $C_4$-invariante. Como $\wp^2$ tem um único polo, de ordem $4$,
em $w=0$, ela define um recobrimento ramificado
$\C/\Lambda\to\mathbb{P}^1$ de grau $4$, donde
$[\C(\wp,\wp'):\C(\wp^2)]=4=|C_4|=[\C(\wp,\wp'):\C(\wp,\wp')^{C_4}]$
e portanto $\C(\wp,\wp')^{C_4}=\C(\wp^2)$.

Finalmente, a inexistência de soluções inteiras não constantes no posto $2$ é o
teorema de Liouville: uma função inteira duplamente periódica é limitada, logo
constante.
\end{proof}

\begin{observacao}
A identidade \eqref{eq:CM} sai da definição: substituindo $\lambda=i\mu$ na
série de $\wp(iw)=\frac1{(iw)^2}+\sum_{\lambda\ne0}
\bigl[\frac1{(iw-\lambda)^2}-\frac1{\lambda^2}\bigr]$ e usando
$i\Lambda=\Lambda$.
\end{observacao}

\subsection*{Censo dos 6144 flexes}

Enumerando exaustivamente $\FF$ com aritmética exata em $\Q(i)$
(Seção~\ref{sec:computacional}) obtemos:

\begin{center}
\renewcommand{\arraystretch}{1.2}
\begin{tabular}{ccrrl}
\toprule
$\operatorname{posto}\Lambda$ & $|P|$ & \#flexes & \% & tipo cristalográfico\\
\midrule
$0$ & $1$ & $4$    & $0{,}07$  & trivial\\
$0$ & $2$ & $36$   & $0{,}59$  & $C_2$\\
$0$ & $4$ & $128$  & $2{,}08$  & $C_4$\\
$1$ & $1$ & $32$   & $0{,}52$  & friso $p1$\\
$1$ & $2$ & $288$  & $4{,}69$  & friso $p2$\\
$2$ & $1$ & $60$   & $0{,}98$  & $p1$\\
$2$ & $2$ & $348$  & $5{,}66$  & $p2$\\
$2$ & $4$ & $5248$ & $85{,}42$ & $p4$\\
\midrule
\multicolumn{2}{l}{\emph{total}} & $6144$ & $100$ & \\
\bottomrule
\end{tabular}
\end{center}

\begin{corolario}[Rigidez]\label{cor:rigidez}
Exatamente $5656$ dos $6144$ flexes de quadrantes ($92{,}06\%$) não admitem
nenhuma função inteira flexionável não constante. Para esses flexes, toda
solução não trivial \emph{tem polos}: o retrato de fase precisa exibir
singularidades.
\end{corolario}

\begin{corolario}\label{cor:trivial}
$\Gamma_\Phi=\{1\}$ se e somente se $A_1=A_2=A_3=A_4$, isto é, se e somente se
$\Phi$ é a restrição a $Q$ de uma rotação global $z\mapsto i^u z$. Há
exatamente $4$ tais flexes, e para eles toda $f$ é flexionável.
\end{corolario}

\begin{corolario}[Multiplicação complexa forçada]\label{cor:cm}
Se $\operatorname{posto}\Lambda_\Phi=2$, então $\Lambda_\Phi$ é um reticulado
contido em $\frac12\Zi$, portanto comensurável com $\Zi$. Consequentemente
$\C/\Lambda_\Phi$ é isógena à curva lemniscática $\C/\Zi$ e
$\End(\C/\Lambda_\Phi)$ é uma ordem em $\Zi$: toda função complexa flexionável
não constante de um flex genérico é uma função elíptica \emph{com multiplicação
complexa por $\Q(i)$}. Se além disso $|P_\Phi|=4$, então $\Lambda_\Phi$ é
quadrado, $g_3=0$ e $j(\C/\Lambda_\Phi)=1728$.
\end{corolario}

\begin{proof}
Proposição~\ref{prop:discreto} dá $\Lambda_\Phi\subset\frac12\Zi$; sendo de
posto $2$, o índice $[\frac12\Zi:\Lambda_\Phi]$ é finito, logo $\Lambda_\Phi$ e
$\Zi$ são comensuráveis e $\Q\Lambda_\Phi=\Q(i)$. Assim
$\End(\C/\Lambda_\Phi)=\{\alpha:\alpha\Lambda_\Phi\subset\Lambda_\Phi\}$ é uma
ordem em $\Q(i)$, estritamente maior que $\Z$. Para $|P|=4$ vale
$i\Lambda_\Phi=\Lambda_\Phi$; então $g_3(i\Lambda)=i^{-6}g_3(\Lambda)=-g_3(\Lambda)$
força $g_3=0$ e $j=1728\,g_2^3/(g_2^3-27g_3^2)=1728$.
\end{proof}

\begin{corolario}[Funções universalmente flexionáveis]\label{cor:universal}
Seja $\Gamma_{\max}=\langle\Gamma_\Phi:\Phi\in\FF\rangle$. Então
$\Gamma_{\max}$ é o grupo cristalográfico $p4$ de reticulado $\Zi$, com centros
de ordem $4$ em $\Zi\cup(\frac{1+i}2+\Zi)$. Em consequência, uma função
meromorfa é $\Phi$-flexionável para \emph{todo} $\Phi\in\FF$ se e somente se
\[
f\in\C\bigl(J_0\bigr),\qquad J_0(z):=\wp\bigl(z;\Zi\bigr)^2 .
\]
\end{corolario}

\section{O que os flexágonos realmente fazem}\label{sec:flexagonos}

\begin{figure}[t]
\centering
\includegraphics[width=\linewidth]{fig/planos.png}
\caption{Os dois planos de papel. À esquerda, o anel de doze quadradinhos do
hexa-tetraflexágono, com as charneiras marcadas por cor (verticais e
horizontais alternam em blocos de três --- é a Proposição~\ref{prop:solta}) e
com moldura nas quatro folhas que carregam a face $1$. À direita, o octomino do
tri-tetraflexágono: das nove adjacências, duas são \emph{cortes} --- o rasgo,
a laranja --- e o que sobra é a tira $a\,b\,c\,f\,g\,h\,e\,d$, colada pelas duas
pontas. A célula $d$, branca dos dois lados, é a aba livre.}\label{fig:planos}
\end{figure}

\begin{figure}[t]
\centering
\includegraphics[width=\linewidth]{fig/tiras.png}
\caption{Os outros dois planos de papel, gerados pelo mesmo código que o modelo
usa (\texttt{hexaflex.tira}). São tiras retas de triângulos equiláteros, com o
último colado de volta no primeiro: $10$ triângulos para o tri-hexaflexágono
($9\cdot2=3\cdot6$) e $19$ para o hexa-hexaflexágono ($18\cdot2=6\cdot6$).
Comparadas com a Figura~\ref{fig:planos}, a diferença que decide tudo não é o
tamanho: é que aqui as charneiras, ao serem levadas para a face, caem nas três
diagonais do hexágono, e essas passam todas pelo centro
(Teorema~\ref{thm:retas}).}\label{fig:tiras}
\end{figure}

\subsection{Os dois modelos de \cite{manual}}
O manual descreve dois tetraflexágonos quadrados. O \emph{hexa-tetraflexágono}
($6$ faces) é montado a partir de um anel de $12$ quadradinhos --- um quadrado
$4\times4$ sem o $2\times2$ central --- \emph{sem cola} \cite[\S1.2]{manual};
a contagem fecha, $12\cdot2=6\cdot4$.

O \emph{tri-tetraflexágono} ($3$ faces) é montado a partir do modelo da
Figura~A.1 de \cite{manual}. Extraindo essa figura pixel a pixel obtém-se, em
cada metade da folha, um \emph{octomino} de $8$ quadradinhos: $6$ com figura e
$2$ brancos --- as abas de cola do roteiro (``cole a face do quadrado branco, a
qual deve estar encarando outro''). Dobrando a folha na linha indicada as duas
metades se superpõem célula a célula; ao todo há $4$ lados brancos em $3$
células, e restam $16-4=12=3\cdot4$ lados com figura. Numerando as células
$a,\dots,e$ na linha de cima e $f,g,h$ na de baixo de uma grade $2\times5$:
\[
\begin{array}{llll}
a:\ E_2/\text{branco} &\quad b:\ E_1/V_1 &\quad c:\ F_2/V_2 &\quad d:\ \text{branco}/\text{branco}\\
e:\ \text{branco}/F_1 &\quad f:\ F_3/E_4 &\quad g:\ V_3/E_3 &\quad h:\ V_4/F_4 .
\end{array}
\]
Qual par de faces brancas é colado não está dito explicitamente, mas fica
determinado por necessidade: enumerando (como abaixo) todas as dobraduras
planas do octomino sem nenhuma restrição, o \emph{único} par de células que
aparece junto em todos os estados que exibem alguma face é $\{a,e\}$ --- as
duas pontas da tira, exatamente como no tri-tetraflexágono clássico. A célula
$d$, branca dos dois lados, é uma aba livre.

\subsection{Mapas de retorno}
Precisamos de uma noção de flex que não dependa de nenhuma convenção de
desenho. Ela é sugerida pelo próprio roteiro de flexão de \cite[p.~6]{manual},
$1 \to 3 \to 4 \to 2 \to 1 \to 5 \to 6 \to 2 \to 1$, em que as faces $1$ e $2$
\emph{reaparecem}.

\begin{definicao}\label{def:retorno}
Um \emph{estado} do flexágono é uma dobradura plana do seu plano sobre um
quadrado $2\times2$. A face $k$ é \emph{exibível} no estado $S$ se os quatro
quadradinhos da face $k$ ocupam, em $S$, as quatro posições distintas do
quadrado, com o lado $k$ para cima; nesse caso $S$ determina uma bijeção $P_S$
das quatro peças nas quatro posições e uma rotação $\rho_S(j)\in\mu_4$ de cada
peça. Se $S$ e $S'$ exibem ambos a face $k$, o \emph{mapa de retorno} é o flex
de quadrantes
\[
\Phi_{S\to S'}=(\sigma,\varepsilon),\qquad \sigma=P_{S'}\circ P_S^{-1},\qquad
\varepsilon_{P_S(j)}=\rho_{S'}(j)\rho_S(j)^{-1}.
\]
\end{definicao}

Se a face $k$ é pintada para ficar correta no estado $S$, no estado $S'$ o
desenho aparece exatamente como $\Phi_{S\to S'}$ o rearranja. A razão
$\rho_{S'}/\rho_S$ elimina qualquer convenção sobre como se desenha o verso do
papel: em ambos os estados o \emph{mesmo} lado de cada quadradinho está para
cima, e um eventual fator de convenção cancela. O mapa de retorno é, portanto,
um invariante do flexágono. Denotamos por $\mathcal{R}$ o conjunto dos mapas de
retorno de um flexágono e por $\Gamma_{\mathcal R}=\langle\Gamma_\Phi:
\Phi\in\mathcal R\rangle$.

\subsection{Enumeração exata das dobraduras planas}
Uma dobradura plana é uma atribuição, a cada quadradinho $C$, de uma isometria
$g_C$ do plano tal que quadradinhos vizinhos $C\sim C'$ de aresta comum $e$
satisfazem $g_{C'}=g_C$ (sem dobra) ou $g_{C'}=g_C\circ\mathrm{ref}_e$ (dobra).
Fixando $g$ numa célula, uma dobradura fica determinada por um bit por
charneira, sujeito ao fechamento dos ciclos independentes do grafo de
adjacência. É essencial percorrer $g_{C_0}\in\{\mathrm{id},\text{reflexão}\}$:
a segunda opção corresponde a virar o flexágono do avesso, e omiti-la elimina
todos os estados em que o verso do papel está para cima. Exigindo por fim que
todas as imagens caiam num bloco $2\times2$ (e, no tri, que $a$ e $e$ fiquem na
mesma posição), obtemos:

\begin{proposicao}\label{prop:enum}
A menos de translação, o anel do hexa-tetraflexágono admite exatamente $70$
dobraduras planas sobre um quadrado $2\times2$, e o octomino colado do
tri-tetraflexágono exatamente $4$. No tri, as três faces são exibíveis, e
apenas a face da frente o é em dois arranjos. No hexa, todas as seis faces são
exibíveis; as faces $1$ e $2$ em $5$ arranjos e as demais em $2$.
\end{proposicao}

A enumeração é um \emph{majorante}: nem toda dobradura plana admite ordenação
de camadas sem auto-interseção, nem é necessariamente alcançável por flexões.
Para o hexa-tetraflexágono, porém, dispomos do \emph{diagrama de flexão} do
próprio projeto \cite{enimar}, que lista os estados efetivamente percorridos.
Nele cada face aparece nos arranjos
\begin{center}
\renewcommand{\arraystretch}{1.2}
\begin{tabular}{cll}
\toprule
faces & arranjos exibidos & mapas de retorno\\
\midrule
$1$, $2$ & $abcd$,\ $badc$,\ $cdab$ & $\sigma_{\mathrm h}$, $\sigma_{\mathrm v}$, $\sigma_{\mathrm h}\sigma_{\mathrm v}$\\
$3$, $6$ & $abcd$,\ $badc$ & $\sigma_{\mathrm h}$\\
$4$, $5$ & $abcd$,\ $cdab$ & $\sigma_{\mathrm v}$\\
\bottomrule
\end{tabular}
\end{center}
onde escrevemos $abcd$ para a palavra que se lê no diagrama (as peças nas
posições superior-esquerda, superior-direita, inferior-esquerda,
inferior-direita), e
\[
\sigma_{\mathrm h}=(12)(34),\qquad \sigma_{\mathrm v}=(14)(23),\qquad
\varepsilon\equiv 1
\]
são, respectivamente, a troca das colunas e a troca das linhas. Esses três
arranjos estão entre os cinco produzidos pela enumeração, o que confirma uma
construção pela outra. Note que \emph{nenhuma peça aparece girada}: os mapas de
retorno dos dois flexágonos são puras permutações de quadrantes --- ``recortar
o quadrado em quatro e trocar as peças de lugar''.

\begin{proposicao}\label{prop:gammas}
Os grupos de discrepância são:
\begin{center}
\renewcommand{\arraystretch}{1.2}
\begin{tabular}{llll}
\toprule
flexágono & face & $\Gamma$ da face & soluções da face\\
\midrule
tri  & frente & $2\Z$ (posto 1) & $h(e^{i\pi z})$\\
tri  & verso, escondida & $\{1\}$ & quaisquer\\
hexa & $3$, $6$ & $2\Z$ (posto 1) & $h(e^{i\pi z})$\\
hexa & $4$, $5$ & $2i\Z$ (posto 1) & $h(e^{\pi z})$\\
hexa & $1$, $2$ & $2\Zi$ (posto 2, $P=1$) & $\C(\wp,\wp')$\\
\midrule
tri  & \emph{todas} & $2\Z$ & $h(e^{i\pi z})$\\
hexa & \emph{todas} & $2\Zi$ & $\C(\wp,\wp')$\\
\bottomrule
\end{tabular}
\end{center}
Em todos os casos $\varepsilon\equiv1$, de modo que $\Gamma_{\mathcal R}$ é um
grupo de \emph{translações} puras: o grupo de pontos é trivial.
\end{proposicao}

\begin{proof}
$\sigma_{\mathrm h}$ com $\varepsilon\equiv1$ tem $A_1(z)=z-1$, $A_2(z)=z+1$,
$A_3(z)=z+1$, $A_4(z)=z-1$, donde $A_2^{-1}A_1=t_{-2}$ e $A_4^{-1}A_1=\mathrm{id}$:
$\Gamma=2\Z$. Analogamente $\sigma_{\mathrm v}$ dá $2i\Z$ e
$\sigma_{\mathrm h}\sigma_{\mathrm v}=(13)(24)$ dá $\langle 2,2i\rangle=2\Zi$.
Os grupos de pontos são triviais porque todos os $\varepsilon_j$ valem $1$.
\end{proof}

\section{Os teoremas para os dois flexágonos}\label{sec:principais}

\begin{teorema}[Tri-tetraflexágono]\label{thm:tri}
Uma função meromorfa é flexionável nas três faces do tri-tetraflexágono se e
somente se ela é $2$-periódica, isto é, se e somente se
\[
\boxed{\;f(z)=h\bigl(e^{i\pi z}\bigr),\qquad h \text{ meromorfa em } \C^\ast
\text{ arbitrária}.\;}
\]
Em particular existem soluções \emph{inteiras} não constantes --- a mais
simples é $f(z)=e^{i\pi z}$ --- e o espaço das soluções é de dimensão infinita.
\end{teorema}

\begin{proof}
Proposição~\ref{prop:gammas} e o caso $(\text{posto }1,\,n=1)$ do
Teorema~\ref{thm:classificacao}, com $\omega=2$.
\end{proof}

\begin{teorema}[Hexa-tetraflexágono]\label{thm:hexa}
Seja $\Lambda=2\Zi$. Uma função meromorfa é flexionável em todas as faces do
hexa-tetraflexágono se e somente se ela é elíptica de reticulado $\Lambda$:
\[
\boxed{\;f\in\C\bigl(\wp(z;2\Zi),\,\wp'(z;2\Zi)\bigr).\;}
\]
\emph{Nenhuma função inteira não constante} serve. Já as faces $3,4,5,6$,
tomadas isoladamente, são satisfeitas por funções inteiras: $e^{i\pi z}$ nas
faces $3$ e $6$, $e^{\pi z}$ nas faces $4$ e $5$. As faces $1$ e $2$ sozinhas
já forçam a elipticidade.

O reticulado é quadrado: $g_3=0$, $j=1728$, $g_2=\varpi^4/4\approx11{,}816982$,
onde $\varpi=2{,}62205\ldots$ é a constante lemniscática. A curva elíptica
subjacente é $y^2=4x^3-g_2x$, com multiplicação complexa por $\Zi$.
\end{teorema}

\begin{proof}
Proposição~\ref{prop:gammas} e o caso $(\text{posto }2,\,n=1)$ do
Teorema~\ref{thm:classificacao}. A não existência de soluções inteiras é
Liouville; $g_3=0$ e $j=1728$ seguem do Corolário~\ref{cor:cm}, e
$g_2(2\Zi)=g_2(\Zi)/16=4\varpi^4/16$.
\end{proof}

\begin{observacao}[a dicotomia]
Vale a pena enunciá-la nitidamente. \emph{As faces que o roteiro de flexão
visita uma única vez são tolerantes; as que ele revisita são rígidas.} No
tri-tetraflexágono só a face da frente reaparece, e ela reaparece ao longo de
\emph{um} eixo: o grupo tem posto $1$ e sobram funções inteiras. No
hexa-tetraflexágono as faces $1$ e $2$ reaparecem ao longo de \emph{dois} eixos
--- é a propriedade que o manual destaca, ``apenas a face de número $1$ pode
ser aberta ao longo de dois eixos diferentes'' --- e duas translações
independentes obrigam a função a ter polos. A rigidez não vem de o flexágono
ser complicado; vem de ele ter uma face que volta por dois caminhos.
\end{observacao}

\begin{observacao}
$\wp(z;2\Zi)$ tem um único polo em $Q=[-1,1]^2$, na origem: a singularidade ---
o ponto onde todas as cores se encontram --- cai exatamente no cruzamento dos
dois cortes, isto é, esconde-se na dobra. Os seus zeros ficam em $\pm1$ e
$\pm i$, nos pontos médios dos lados.
\end{observacao}

\section{Fora do mundo holomorfo não há rigidez}\label{sec:naoanalitico}
Vale a pena isolar onde a rigidez foi usada: apenas no princípio da identidade.

\begin{proposicao}
Seja $\Sigma\subset Q$ a união dos dois cortes (a ``cruz'' $\{xy=0\}\cap Q$) e
$E_\Phi:=\bigcup_{j}A_j^{-1}\bigl(\Sigma\cap\partial Q_{\sigma(j)}\bigr)$, um
subconjunto de dimensão $1$ de $Q$. Se exigimos apenas que $\Phi_*f$ seja
\emph{contínua} em $Q$ (com $f$ contínua), a condição envolve somente
$f|_{E_\Phi}$. Em particular, qualquer $f$ contínua que se anule numa vizinhança
de $E_\Phi$ é flexionável nesse sentido, e o espaço das soluções é de dimensão
infinita para todo $\Phi\in\FF$.
\end{proposicao}

O mesmo vale na categoria real-analítica: duas funções real-analíticas em
$\R^2$ podem coincidir num segmento sem coincidir (por exemplo $0$ e $y$ ao
longo de $y=0$), de modo que a colagem ao longo dos cortes \emph{não} propaga
informação. É a rigidez do princípio da identidade complexo --- o fato de que
uma função holomorfa é determinada por sua restrição a um segmento --- que
transforma quatro condições de fronteira em uma equação funcional global, e daí
em funções elípticas. Nesse sentido, a resposta ao problema do título é um
fenômeno estritamente da análise complexa.

\section{Verificação computacional}\label{sec:computacional}

Todos os enunciados foram verificados por programas independentes.

\begin{enumerate}
\item \textbf{Simulação da dobradura} (\texttt{flexagon\_sim.py}). Aritmética
exata em $\Q(i)$; reproduz os cinco diagramas de montagem de
\cite[p.~5]{manual}, inclusive o passo final de ``desenrolar e enrolar ao
contrário'', e devolve o estado montado com todas as posições e orientações.

\item \textbf{Extração do modelo impresso} (\texttt{plano\_tri.py}). A
Figura~A.1 de \cite{manual} é lida pixel a pixel; confirma-se que cada metade é
um octomino, que as duas metades se superpõem célula a célula sob a dobra
indicada, que sobram $12$ lados com figura, e que o gerador de planos do projeto
reproduz exatamente esse modelo.

\item \textbf{Enumeração das dobraduras} (\texttt{foldings.py},
\texttt{foldings\_tri.py}, \texttt{mecanica.py}). Aritmética inteira exata em
$\Zi$ (coordenadas dobradas); produz as $70$ dobraduras planas do hexa e as $4$
do tri, os arranjos exibíveis de cada face e os mapas de retorno da
Proposição~\ref{prop:enum}. O estado montado devolvido pelo item 1 é
reencontrado dentro da lista, e os arranjos do diagrama de \cite{enimar} estão
entre os enumerados --- três construções independentes que concordam.

\item \textbf{Reticulados} (\texttt{reticulados.py}). Percorre os $5656$ flexes
de posto $2$, reduz cada reticulado por Gauss, testa se é da forma $\alpha\Zi$
e calcula $j(\tau)$ por séries de Eisenstein: são os dados do
Teorema~\ref{thm:reticulados}.

\item \textbf{Censo de $\FF$} (\texttt{quadflex.py}). Os $6144$ flexes são
enumerados e, para cada um, $\Gamma_\Phi$ é calculado exatamente: o grupo de
pontos por fecho em $\mu_4$ e o reticulado por geradores de Schreier seguidos
de forma normal de Hermite sobre $\Z$. Isso produz a tabela da
Seção~\ref{sec:classificacao} e os Corolários
\ref{cor:rigidez}--\ref{cor:universal}. O Lema~\ref{lem:rho} foi testado em
$400$ flexes aleatórios contra as $4$ rotações globais.

\item \textbf{Modelo com camadas} (\texttt{verify\_camadas.py},
\texttt{verify\_geral.py}). $65$ testes exatos, em aritmética inteira:
a fórmula de paridade e a fórmula dos intervalos contra a enumeração dos
$216\,768$ estados; a regra dos extremos e a sua versão espelhada; a
decomposição completa em $20\,405$ componentes (soma $=216\,768$, ilha $=19\,200$,
maior grão $=200$); $W^+(D)$ nos dois reticulados; a flexão de pinça; o
tri-hexaflexágono com $\Gamma=C_3$; a roda de dois protótipos com $\Gamma=C_{12}$;
e os quatro catálogos --- não só a existência dos arquivos, mas que os
quatro \texttt{catalogo.pdf} são PDFs de verdade e que nenhuma fórmula
impressa traz sinal duplo.

\item \textbf{Testes funcionais} (\texttt{verify2.py}).
$\wp$ e $\wp'$ são avaliadas por funções teta de Jacobi com base de Gauss
reduzida e conferidas contra a soma de Eisenstein direta e contra diferenças
finitas (concordância $\sim10^{-6}$ e $\sim10^{-7}$). Testamos
$\max_j\|f\circ A_1^{-1}-f\circ A_j^{-1}\|$ em $1500$ pontos aleatórios de $Q$
contra todos os mapas de retorno:
\begin{center}
\renewcommand{\arraystretch}{1.15}
\begin{tabular}{llc}
\toprule
$f$ & mapas de retorno & resíduo relativo máximo\\
\midrule
$\wp(z;2\Zi)$                 & todos os do hexa   & $7{,}2\times10^{-15}$\\
$R(\wp)+S(\wp)\wp'$           & todos os do hexa   & $\le 2{,}2\times10^{-13}$\\
$e^{i\pi z}$                  & faces $3$, $6$; todo o tri & $9{,}1\times10^{-16}$\\
$e^{\pi z}$                   & faces $4$, $5$     & $9{,}4\times10^{-16}$\\
\midrule
$e^{i\pi z}$                  & face $1$ (diagonal) & $2{,}5\times10^{2}$ \ (falha)\\
$e^{z}$, $z^2$                & todos              & $\ge 3{,}9$ \ (falha)\\
\bottomrule
\end{tabular}
\end{center}
Verificamos também $\wp(iz)=-\wp(z)$ (resíduo $3\times10^{-15}$), $g_3=0$ e
$j=1728$ para $\Zi$, $2\Zi$ e $(1+i)\Zi$.

\item \textbf{Geração e desenho} (\texttt{gerador.py}). Ver
Seção~\ref{sec:gerador}: cada imagem produzida é testada antes de ser salva.
\end{enumerate}

\begin{figure}[t]
\centering
\begin{subfigure}[b]{0.235\textwidth}
  \includegraphics[width=\linewidth]{fig/n_wp_diag_antes.png}
  \caption{$f=\wp(z;2\Zi)$}
\end{subfigure}\hfill
\begin{subfigure}[b]{0.235\textwidth}
  \includegraphics[width=\linewidth]{fig/n_wp_diag_depois.png}
  \caption{$\Phi_*f$: costura invisível}
\end{subfigure}\hfill
\begin{subfigure}[b]{0.235\textwidth}
  \includegraphics[width=\linewidth]{fig/n_exp_diag_antes.png}
  \caption{$f=e^{2z}$}
\end{subfigure}\hfill
\begin{subfigure}[b]{0.235\textwidth}
  \includegraphics[width=\linewidth]{fig/n_exp_diag_depois.png}
  \caption{$\Phi_*f$: quebra}
\end{subfigure}
\caption{O mapa de retorno diagonal da face $1$ ($\sigma=(13)(24)$,
$\varepsilon\equiv1$: os quadrantes opostos trocam de lugar, sem girar)
aplicado a uma solução e a uma não-solução. Em (b) a figura é literalmente a
mesma de (a). As linhas brancas marcam os cortes.}
\end{figure}

\begin{figure}[t]
\centering
\begin{subfigure}[b]{0.235\textwidth}
  \includegraphics[width=\linewidth]{fig/n_exp_horz_antes.png}
  \caption{$f=e^{i\pi z}$}
\end{subfigure}\hfill
\begin{subfigure}[b]{0.235\textwidth}
  \includegraphics[width=\linewidth]{fig/n_exp_horz_depois.png}
  \caption{faces $3$, $6$: cola}
\end{subfigure}\hfill
\begin{subfigure}[b]{0.235\textwidth}
  \includegraphics[width=\linewidth]{fig/n_exp_ihdiag_antes.png}
  \caption{a mesma $f$}
\end{subfigure}\hfill
\begin{subfigure}[b]{0.235\textwidth}
  \includegraphics[width=\linewidth]{fig/n_exp_ihdiag_depois.png}
  \caption{face $1$: quebra}
\end{subfigure}
\caption{A solução \emph{inteira} $e^{i\pi z}$ das faces $3$ e $6$ do
hexa-tetraflexágono (Proposição~\ref{prop:gammas}) --- a mesma que resolve
\emph{todo} o tri-tetraflexágono (Teorema~\ref{thm:tri}). Ela não resolve a
face~$1$ do hexa: nenhuma função inteira não constante resolve todas.}
\end{figure}

\section{Gerando a família completa}\label{sec:gerador}

A classificação do Teorema~\ref{thm:classificacao} é construtiva, e isso permite
mais do que exibir exemplos: permite \emph{parametrizar} o conjunto das
soluções. Em cada uma das oito classes a órbifold $\C/\Gamma_\Phi$ é uma
superfície de Riemann com uma função geradora $J$ --- um Hauptmodul --- tal que
toda solução meromorfa é $f=R(J)$ com $R$ racional (no caso
$(\text{posto }2,\,|P|=1)$, $f=R_1(\wp)+R_2(\wp)\wp'$):
\begin{center}
\renewcommand{\arraystretch}{1.25}
\begin{tabular}{ccl}
\toprule
$\operatorname{posto}\Lambda$ & $n=|P|$ & gerador $J$ da órbifold\\
\midrule
$0$ & $1$ & $J=z$ (toda $f$ serve)\\
$0$ & $n$ & $J=(z-p)^n$\\
$1$ & $1$ & $J=e^{2\pi i(z-p)/\omega}$\\
$1$ & $2$ & $J=\cos\bigl(2\pi (z-p)/\omega\bigr)$\\
$2$ & $1$ & $(\wp,\wp')$ (dois geradores)\\
$2$ & $2$ & $J=\wp(z-p)$\\
$2$ & $4$ & $J=\wp(z-p)^2$\\
\bottomrule
\end{tabular}
\end{center}

O programa \texttt{gerador.py} implementa exatamente isso. Dado um flexágono ou
um flex arbitrário, ele calcula $\Gamma$ com aritmética exata em $\Q(i)$,
identifica a classe, constrói $J$, sorteia $R$ e --- este é o ponto --- confere
numericamente que a $f$ resultante é de fato flexionável antes de desenhá-la. As
raízes e os polos de $R$ são sorteados entre os \emph{valores} que $J$
efetivamente assume no quadrado, o que garante que os zeros e polos de $f$ caiam
dentro da figura; sem esse cuidado a maioria dos retratos sai quase constante.
Uma pequena varredura de qualidade (cobertura do círculo cromático) descarta o
resto. Na prática todas as tentativas passam no teste, com resíduo da ordem de
$10^{-14}$.

\begin{figure}[t]
\centering
\begin{subfigure}[b]{0.30\textwidth}
  \includegraphics[width=\linewidth]{fig/hexa_face1.png}
\end{subfigure}\hfill
\begin{subfigure}[b]{0.30\textwidth}
  \includegraphics[width=\linewidth]{fig/hexa_face3.png}
\end{subfigure}\hfill
\begin{subfigure}[b]{0.30\textwidth}
  \includegraphics[width=\linewidth]{fig/hexa_face5.png}
\end{subfigure}\\[4pt]
\begin{subfigure}[b]{0.30\textwidth}
  \includegraphics[width=\linewidth]{fig/hexa_face2.png}
\end{subfigure}\hfill
\begin{subfigure}[b]{0.30\textwidth}
  \includegraphics[width=\linewidth]{fig/hexa_face4.png}
\end{subfigure}\hfill
\begin{subfigure}[b]{0.30\textwidth}
  \includegraphics[width=\linewidth]{fig/hexa_face6.png}
\end{subfigure}
\caption{Seis faces geradas automaticamente para o hexa-tetraflexágono:
$f=R_k(\wp(z;2\Zi))$ com racionais $R_k$ sorteadas. São seis desenhos
diferentes e todos sobrevivem a todos os mapas de retorno (resíduos
$\le 1{,}2\times10^{-13}$).}
\end{figure}
\section{A teoria geral: faces ladrilhadas}\label{sec:geral}

Tudo o que foi feito até aqui para o quadrado vale, sem alteração, para
qualquer face ladrilhada por polígonos congruentes. Vale a pena enunciá-lo
assim, porque é o que faz os flexágonos triangulares (e quaisquer outros)
caírem no mesmo teorema.

\begin{definicao}
Seja $\mathcal T$ uma tesselação do plano por polígonos congruentes, com
reticulado de translações $L$ e grupo de rotações $C_n$; pela restrição
cristalográfica $n\in\{1,2,3,4,6\}$. Uma \emph{face} é uma união finita
$D=T_1\cup\dots\cup T_m$ de ladrilhos. Um \emph{flex} de $D$ é uma bijeção
$\Phi:D\to D$, definida fora das arestas, que é isometria direta em cada
$\mathring T_j$ e permuta os ladrilhos:
\[
\Phi|_{\mathring T_j}=A_j,\qquad A_j(z)=\varepsilon_j(z-c_j)+c_{\sigma(j)},
\]
onde $c_j$ é o baricentro de $T_j$ e $\varepsilon_j$ percorre as rotações que
levam a forma de $T_j$ na de $T_{\sigma(j)}$ --- uma classe lateral do grupo de
rotações do ladrilho. O grupo dos flexes tem ordem $m!\cdot|\mathrm{Rot}(T)|^m$.
\end{definicao}

Para o quadrado $2\times2$ de quadradinhos: $n=4$, $m=4$,
$\mathrm{Rot}(T)=C_4$, $|\FF|=4!\cdot4^4=6144$. Para o hexágono de seis
triângulos equiláteros: $n=6$, $m=6$, $\mathrm{Rot}(T)=C_3$ (o triângulo só
admite rotações de $120^\circ$ em torno do seu baricentro) e
$|\FF_\triangle|=6!\cdot3^6=524\,880$.

\begin{teorema}[Teorema geral]\label{thm:geral}
O Teorema~\ref{thm:principal} vale literalmente nesse contexto: $f$ é
$\Phi$-flexionável se e somente se $f$ é invariante por
$\Gamma_\Phi=\langle A_k^{-1}A_j\rangle$. Além disso
\[
\Gamma_\Phi\ \le\ C_n\ltimes\tfrac1d\,\mathcal O_n,\qquad
\mathcal O_4=\Zi,\quad \mathcal O_3=\mathcal O_6=\Z[\zeta_3],
\]
onde $d$ é o denominador dos baricentros ($d=2$ no caso quadrado, $d=3$ no
triangular). Em particular $\Gamma_\Phi$ é discreto, e se o seu reticulado de
translações $\Lambda$ tem posto $2$ então $\Lambda$ é comensurável com
$\mathcal O_n$; se ainda $|P_\Phi|\ge3$, então $\Lambda$ é um ideal de
$\mathcal O_n$ e $\C/\Lambda$ tem multiplicação complexa:
\[
n=4:\quad j=1728 \ \text{(lemniscática)};\qquad
n=3,6:\quad j=0 \ \text{(equianarmônica)}.
\]
\end{teorema}

\begin{proof}
A demonstração do Teorema~\ref{thm:principal} usa apenas que os ladrilhos são
abertos conexos, que $D$ tem interior conexo e o princípio da identidade;
nada de específico ao quadrado. Para a segunda parte: $d\,c_j\in\mathcal O_n$ e
$\varepsilon_j\in\mu_n\subset\mathcal O_n^\times$, logo
$d\beta_j=d(c_{\sigma(j)}-\varepsilon_jc_j)\in\mathcal O_n$ e todo
$A_k^{-1}A_j$ tem a forma $z\mapsto\zeta z+t$ com $\zeta\in\mu_n$ e
$t\in\frac1d\mathcal O_n$; esse conjunto é um grupo discreto. Se $|P|\ge3$
então $\zeta_n\Lambda=\Lambda$ com $n\ge3$, isto é, $\Lambda$ é um
$\mathcal O_n$-submódulo de $\Q(\zeta_n)$ finitamente gerado, ou seja um ideal
fracionário; como $\mathcal O_n$ é principal, $\Lambda=\alpha\mathcal O_n$ e a
homotetia $z\mapsto z/\alpha$ leva $\Lambda$ em $\mathcal O_n$. Os invariantes
de $\C/\Zi$ e $\C/\Z[\zeta_3]$ são $j=1728$ e $j=0$.
\end{proof}

\begin{teorema}[Classificação universal]\label{thm:universal}
Seja $\Phi$ um flex de uma face ladrilhada e $f$ meromorfa não constante.
Então $f$ é $\Phi$-flexionável se e somente se $f$ pertence ao corpo indicado,
onde $n=|P_\Phi|\in\{1,2,3,4,6\}$, $\omega$ gera $\Lambda$ no posto $1$,
$\wp=\wp(\,\cdot\,;\Lambda)$ e $p$ é um centro de rotação:
\begin{center}
\renewcommand{\arraystretch}{1.3}
\begin{tabular}{ccll}
\toprule
$\operatorname{posto}\Lambda$ & $|P|$ & gerador da órbifold $\C/\Gamma_\Phi$ & inteiras?\\
\midrule
$0$ & $1$ & --- (toda $f$ serve) & sim\\
$0$ & $n$ & $(z-p)^n$ & sim\\
$1$ & $1$ & $e^{2\pi i(z-p)/\omega}$ & sim\\
$1$ & $2$ & $\cos\bigl(2\pi(z-p)/\omega\bigr)$ & sim\\
$2$ & $1$ & $\wp,\ \wp'$ \quad(dois geradores) & \textbf{não}\\
$2$ & $2$ & $\wp(z-p)$ & \textbf{não}\\
$2$ & $3$ & $\wp'(z-p)$;\ $\Lambda$ hexagonal, $g_2=0$, $j=0$ & \textbf{não}\\
$2$ & $4$ & $\wp(z-p)^2$;\ $\Lambda$ quadrado, $g_3=0$, $j=1728$ & \textbf{não}\\
$2$ & $6$ & $\wp'(z-p)^2$;\ $\Lambda$ hexagonal, $g_2=0$, $j=0$ & \textbf{não}\\
\bottomrule
\end{tabular}
\end{center}
\end{teorema}

\begin{figure}[p]
\centering
\includegraphics[width=\linewidth]{fig/atlas.png}
\caption{As nove classes do Teorema~\ref{thm:universal}, cada uma pelo retrato
de fase do seu gerador $J$ (com $\omega=2$ e $p=0$; a cor é o argumento, as
faixas escuras são curvas de nível de $|J|$). A tabela lê-se na figura. No
posto $0$ o gerador é $(z-p)^n$ e tem no centro um zero de ordem $n=|P|$: uma
volta de cor quando $|P|=1$, três quando $n=3$. No posto $2$ o centro é um
\emph{polo}, de ordem $|P|$ --- dois em $\wp$, três em $\wp'$, quatro em
$\wp^2$, seis em $\wp'^2$ ---, com a única exceção de $(2,1)$, onde nenhum
gerador sozinho serve e desenhamos $\wp'$. No posto $1$ não há ponto
distinguido: $e^{2\pi i(z-p)/\omega}$ nunca se anula e os dois zeros do
cosseno ficam fora do centro. O posto de $\Lambda$ é o número de direções em
que o desenho se repete, e a simetria do retículo de polos separa $\Lambda$
quadrado ($j=1728$) de $\Lambda$ hexagonal ($j=0$). Nas três primeiras linhas
há soluções inteiras não constantes; nas seis de posto $2$, não.}
\label{fig:atlas}
\end{figure}

\begin{proof}
Como antes, trata-se de descrever as funções meromorfas na órbifold
$\C/\Gamma_\Phi$; os casos de posto $\le1$ e os de posto $2$ com $|P|\le2$ já
foram feitos. Para $|P|=3$: $\zeta_3\Lambda=\Lambda$ e
$\wp(\zeta_3w)=\zeta_3^{-2}\wp(w)=\zeta_3\wp(w)$, donde
$\wp'(\zeta_3w)=\zeta_3^{-3}\wp'(w)=\wp'(w)$: $\wp'$ é $C_3$-invariante e tem
grau $3=|C_3|$, logo gera o corpo de invariantes. Para $|P|=6$, $C_6=\langle
-1\rangle\times\langle\zeta_3\rangle$ e $\wp'$ é ímpar, de modo que $\wp'^2$ é
$C_6$-invariante de grau $6=|C_6|$. Para $|P|=4$, $\wp(iw)=-\wp(w)$ e $\wp^2$
tem grau $4$. Em cada caso o índice bate com a ordem do grupo, e o corpo de
invariantes é gerado por esse elemento.
\end{proof}

\subsection*{Os dois censos}
Enumerando exaustivamente os dois grupos de flexes com aritmética exata em
$\Zi$ e em $\Z[\zeta_3]$ (Seção~\ref{sec:computacional}):
\begin{center}
\renewcommand{\arraystretch}{1.15}
\begin{tabular}{ccrrrr}
\toprule
& & \multicolumn{2}{c}{quadrado ($|\FF|=6144$)} & \multicolumn{2}{c}{triangular ($|\FF|=524\,880$)}\\
posto & $|P|$ & \#flexes & \% & \#flexes & \%\\
\midrule
$0$&$1$&    $4$&$0{,}07$&      $6$&$0{,}00$\\
$0$&$2$&   $36$&$0{,}59$&     $54$&$0{,}01$\\
$0$&$3$&    ---&---&         $246$&$0{,}05$\\
$0$&$4$&  $128$&$2{,}08$&      ---&---\\
$0$&$6$&    ---&---&         $738$&$0{,}14$\\
$1$&$1$&   $32$&$0{,}52$&     $54$&$0{,}01$\\
$1$&$2$&  $288$&$4{,}69$&    $522$&$0{,}10$\\
$2$&$1$&   $60$&$0{,}98$&    $156$&$0{,}03$\\
$2$&$2$&  $348$&$5{,}66$&  $1\,368$&$0{,}26$\\
$2$&$3$&    ---&---&      $52\,026$&$9{,}91$\\
$2$&$4$&$5\,248$&$85{,}42$&    ---&---\\
$2$&$6$&    ---&---&     $469\,710$&$89{,}49$\\
\midrule
\multicolumn{2}{l}{\emph{admitem inteiras}} & $488$ & $7{,}94$ & $1\,620$ & $0{,}31$\\
\multicolumn{2}{l}{\emph{forçam polos}} & $5\,656$ & $92{,}06$ & $523\,260$ & $99{,}69$\\
\bottomrule
\end{tabular}
\end{center}

\begin{corolario}[Rigidez, forma geral]\label{cor:rigidez2}
Um flex de quadrantes tomado ao acaso não admite função inteira não constante
com probabilidade $92{,}06\%$; um flex triangular, com probabilidade
$99{,}69\%$. E em $85{,}4\%$ (resp. $99{,}4\%$) dos casos a curva elíptica
subjacente é \emph{forçada}: lemniscática ($j=1728$) no caso quadrado,
equianarmônica ($j=0$) no triangular.
\end{corolario}

\section{O modelo com camadas e o teorema das retas de dobra}\label{sec:camadas}

Até aqui um estado do brinquedo foi uma \emph{dobradura plana}: uma isometria
$g_C$ por célula. Isso ignora duas coisas que o papel não ignora --- a ordem
em que as células se empilham, e o facto de só se ver a de cima. Repô-las
decide a pergunta que ficou em aberto e, mais do que isso, resolve de uma vez
todos os flexágonos.

\begin{definicao}[estado dobrado]
Um \emph{estado dobrado} do plano de um flexágono é um par $(g,\pi)$: uma
isometria $g_C$ por célula e, em cada posição $p$, uma ordem total $\pi_p$ (de
baixo para cima) das células que aí caem, sujeita às três condições de
não-atravessamento de Justin --- taco--taco, taco--tortilha e
tortilha--tortilha. Elas são necessárias e suficientes para que a ordem
descreva papel sem auto-intersecção \cite{demaine}. A face $k$ está
\emph{exibida} quando as suas células são as de \emph{topo} das quatro
posições, com o lado $k$ para cima: é o que se vê, e é sobre o que se vê que a
função está pintada.
\end{definicao}

Uma reta $\ell$ é \emph{admissível} para a face $D$ quando $\ell\cap D$ é uma
união de arestas de ladrilhos --- dobrar não pode vincar um ladrilho --- e a
reflexão $r_\ell$ leva ladrilhos de $\mathcal T$ em ladrilhos de $\mathcal T$.
Um \emph{flex} é então: dobrar ao longo de uma reta admissível (a metade móvel
passa por cima ou por baixo da outra) e reabrir ao longo de outra reta
admissível, fazendo passar para o outro lado um \emph{bloco contíguo de
camadas} do topo ou do fundo --- é o que o dedo faz ao abrir um livro. Só
podem ser cortadas charneiras que assentem sobre a reta de abertura, e o
resultado tem de ser de novo um estado dobrado válido.

\begin{lema}[retas admissíveis]\label{lem:retas}
Para o quadrado $2\times2$ de quadradinhos as retas admissíveis são as duas
medianas e os quatro lados: seis retas. Para o hexágono de seis triângulos
equiláteros são as três diagonais longas --- e \emph{todas passam pelo centro}
$O$.
\end{lema}

\begin{proof}
No quadrado, as arestas do ladrilhamento jazem em $x\in\{0,\pm1\}$ e
$y\in\{0,\pm1\}$; as seis retas são admissíveis, e dobrar por um lado desliza
o quadrado um ladrilho --- é o passo de abertura do flex. No hexágono, as
arestas são os seis raios e os seis lados; raios opostos alinham-se dois a
dois nas três diagonais longas, ao passo que dois lados consecutivos fazem
$120^\circ$ e nunca se alinham, de modo que uma reta admissível contém no
máximo um lado. Dobrar por esse lado levaria a folha incidente para fora do
hexágono, e a união dos ladrilhos restantes com a sua imagem não é um hexágono
de seis triângulos em torno de um vértice comum (a imagem não toca $O$): o
contorno deixaria de ser a face. Restam as três diagonais.
\end{proof}

\begin{figure}[t]
\centering
\includegraphics[width=\linewidth]{fig/retas_dobra.png}
\caption{As retas admissíveis das duas faces. No quadrado há seis --- as duas
medianas e os quatro lados --- e duas paralelas compõem-se numa translação: é
daí que vem o posto $2$. No hexágono há três, as diagonais longas, e
\emph{todas passam pelo centro}: os lados não servem porque levariam papel para
fora da face. Toda a diferença entre $\C(\wp,\wp')$ e $h\big((z-O)^3\big)$ está
nesta figura.}\label{fig:retas}
\end{figure}

\begin{teorema}[teorema das retas de dobra]\label{thm:retas}
Seja $W(D)$ o grupo gerado pelas reflexões nas retas admissíveis de $D$ e
$W^+(D)$ o seu subgrupo de isometrias diretas. Se dois estados de um flexágono
ligados por flexes exibem a mesma face, então o mapa de retorno $\Phi$ entre
eles satisfaz $A_j\in W^+(D)$ para todo $j$. Em particular
$\Gamma_\Phi\le W^+(D)$.
\end{teorema}

\begin{proof}
Cada dobra e cada abertura substitui $g_C$ por $r_\ell\circ g_C$ nas células
que se movem, com $\ell$ admissível. Ao longo de uma sequência de flexes cada
folha $L$ tem portanto $g'_L=h_L\circ g_L$ com $h_L\in W(D)$. Se a face está
exibida nos dois estados, $L$ mostra o mesmo lado para cima nos dois, logo
$h_L$ é direta: $h_L\in W^+(D)$. Pela definição de mapa de retorno,
$A_j=g'_L\circ g_L^{-1}=h_L$ para a folha $L$ que ocupa o ladrilho $j$, e
$\Gamma_\Phi=\langle A_k^{-1}A_j\rangle\le W^+(D)$.
\end{proof}

\begin{corolario}[caso quadrado]\label{cor:quad}
$W$ é gerado pelas reflexões em $x,y\in\{0,\pm1\}$, logo $W^+$ é o grupo $p2$
das translações de $2\Zi$ e das meias-voltas em torno dos pontos de $\Zi$.
Assim $\Lambda\subseteq2\Zi$ e $|P|\le2$: das nove classes do
Teorema~\ref{thm:universal} só $p1$ e $p2$ (e as degeneradas de posto menor)
podem ocorrer num tetraflexágono.
\end{corolario}

\begin{corolario}[caso hexagonal]\label{cor:hex}
As três diagonais fazem $60^\circ$ entre si, logo $W$ é o grupo diedral de
ordem $6$ que fixa $O$ e $W^+=C_3$. Assim $\Gamma$ é um grupo cíclico
\emph{finito} de rotações em torno de $O$, de ordem $m\mid3$, e
\[
\{f \text{ meromorfa}:\ f\circ\Phi^{-1}\text{ meromorfa para todo flex}\}
=\C\big((z-O)^m\big),\qquad\text{isto é}\qquad f(z)=h\big((z-O)^m\big).
\]
Em particular \emph{nenhum} hexaflexágono força funções elípticas, todos
admitem soluções inteiras não constantes --- por exemplo $(z-O)^m$ --- e
\emph{cada face de qualquer hexaflexágono aparece em no máximo três
arranjos}. Este último enunciado é exatamente o que Gardner observou
experimentalmente em $1956$ (Teorema~\ref{thm:hexaflex} e a discussão que o
segue): o limite $3$ é um teorema, não um acidente do hexa-hexa.
\end{corolario}

\begin{observacao}[$W^+(D)$ calcula-se face a face]\label{obs:wgrupo}
O cálculo acima é mecânico e está implementado em \texttt{wgrupo.py} para
polióminos e poliamantes: enumeram-se as retas do reticulado que contêm
arestas da face e forma-se o grupo gerado pelos produtos de duas reflexões.
Distinguem-se duas colunas --- \emph{confinada} (só as retas interiores a $D$:
o brinquedo nunca se abre para fora da face) e \emph{livre} (todas as retas
que tocam $D$: ele pode abrir-se por um lado e deslizar) ---, e $\Gamma$ fica
sempre entre as duas.
\begin{center}
\renewcommand{\arraystretch}{1.15}
\begin{tabular}{lll}
\toprule
face & $W^+$ confinada & $W^+$ livre\\
\midrule
quadrado $2\times2$ (tetraflexágonos) & $C_2$ & posto $2$, $|P|=2$\\
dominó $1\times2$ & trivial & posto $2$, $|P|=2$\\
tira $1\times3$ & posto $1$, $|P|=1$ & posto $2$, $|P|=2$\\
retângulo $2\times3$ & posto $1$, $|P|=2$ & posto $2$, $|P|=2$\\
quadrado $3\times3$ & posto $2$, $|P|=2$ & posto $2$, $|P|=2$\\
\midrule
hexágono de $6$ triângulos (hexaflexágonos) & $C_3$ & posto $2$, $|P|=3$\\
losango ($2$ triângulos) & trivial & posto $2$, $|P|=3$\\
trapézio ($3$ triângulos) & $C_3$ & posto $2$, $|P|=3$\\
\bottomrule
\end{tabular}
\end{center}
\end{observacao}

\noindent O contraste entre os dois corolários tem uma causa só, e é
geométrica: no quadrado há retas admissíveis que \emph{não} passam pelo centro
(os quatro lados), e é a composição de duas reflexões paralelas que produz as
translações e, com elas, o posto $2$; no hexágono todas passam pelo centro, e
o grupo colapsa para rotações.

\subsection*{O hexa-tetraflexágono, decidido}

O Corolário~\ref{cor:quad} deixa exatamente duas possibilidades, $|P|=1$ e
$|P|=2$; é a busca exaustiva no modelo com camadas que escolhe. Implementámos
estados dobrados, exibição e flex em \texttt{flexcamadas.py} e percorremos a
órbita do estado montado, semeando-a com \emph{todas} as ordens de camadas
compatíveis com esse estado --- o que torna a órbita um majorante da
alcançabilidade física e as exclusões incondicionais.

\begin{teorema}\label{thm:decidido}
O anel de $12$ quadradinhos admite $216\,768$ estados dobrados válidos sobre o
quadrado $2\times2$ (num referencial fixo); a órbita do estado montado pelo
flex tem $19\,200$. Nessa órbita as faces $1$ e $2$ exibem-se em exatamente
três arranjos e as faces $3,\dots,6$ em exatamente dois --- precisamente o
diagrama de flexão de \cite{enimar}. Todos os mapas de retorno têm
$\varepsilon$ constante, portanto
\[
\Gamma=2\Zi,\qquad |P|=1,\qquad
\boxed{\ \text{soluções}=\C(\wp,\wp')\ \text{do reticulado quadrado }2\Zi\ }
\]
--- todas as funções elípticas de $2\Zi$, e não apenas as pares.
\end{teorema}

A comparação é instrutiva. O modelo puramente combinatório --- todos os
estados dobrados, sem alcançabilidade --- dá cinco arranjos para as faces $1$
e $2$ e $|P|=2$, isto é $\C(\wp)$. O majorante em que a abertura pode levar
para o outro lado \emph{qualquer} conjunto admissível de camadas, e não um
bloco contíguo, dá também cinco e $|P|=2$. A diferença entre três e cinco é
exatamente a diferença entre abrir um livro por uma página e separar folhas
não contíguas: a primeira é um flex, a segunda desmonta o brinquedo. E os dois
arranjos que só o majorante alcança são precisamente os únicos com
$\varepsilon$ não constante --- os únicos, portanto, capazes de pôr uma
meia-volta em $\Gamma$. A órbita é fechada pelas meias-voltas e pelos dois
espelhos, mas \emph{não} pelas rotações de um quarto de volta: um
tetraflexágono nunca volta rodado $90^\circ$.

\subsection*{Os outros três}

O plano do tri-tetraflexágono precisa de uma correção. Das nove adjacências do
octomino da Figura~A.1 de \cite{manual}, duas são \emph{cortes} --- o rasgo do
tri-tetraflexágono clássico ---, e o que sobra é a tira
$a\,b\,c\,f\,g\,h\,e\,d$ com as duas pontas coladas. Tira e colagem ficam
determinadas por necessidade (\texttt{busca\_tri.py}): das três tiras
hamiltonianas possíveis e de todas as colagens branco-com-branco, só essa
exibe as \emph{três} faces no topo das pilhas, e só ela reproduz o diagrama
de \cite{enimar} ($F$ com dois arranjos, $V$ e $E$ com um). O grupo não muda:
$\Gamma=2\Z$, posto $1$, soluções $h(e^{i\pi z})$.

Nos dois hexaflexágonos o Corolário~\ref{cor:hex} faz quase todo o trabalho:
um mapa de retorno é uma rotação em torno de $O$, o que em coordenadas de
ladrilho se lê $\sigma(j)\equiv j+\varepsilon_j\pmod 6$. Filtrando o censo
combinatório por essa condição (\texttt{rotflex.py}):

\begin{center}
\renewcommand{\arraystretch}{1.15}
\begin{tabular}{lrrl}
\toprule
& mapas combinatórios & rotações em torno de $O$ & $\Gamma$\\
\midrule
tri-hexaflexágono  & $4$   & $0$  & --- \\
hexa-hexaflexágono & $107$ & $35$ & $C_3$\\
\bottomrule
\end{tabular}
\end{center}

(No tri-hexaflexágono o censo combinatório fixa o referencial --- \texttt{hexaflex.enumerar}
percorre só $g_0\in\{\mathrm{ID},\mathrm{REF}\}$ --- e por isso não vê as
rotações que a flexão realiza; é o modelo com camadas, com a pinça, que as
encontra.)

\begin{teorema}\label{thm:hexaflex}
No hexa-hexaflexágono $\Gamma=C_3$ --- o máximo permitido pelo
Corolário~\ref{cor:hex} ---, gerado pela rotação de $120^\circ$ em torno do
centro, e as soluções são exatamente $f(z)=h\big((z-O)^3\big)$ com $h$
meromorfa, entre elas $(z-O)^3$, inteira e não constante. E o mesmo vale para o
tri-hexaflexágono: com a flexão de pinça, a órbita de um exemplar montado tem
$660$ estados, as suas três faces --- as seis folhas de cima, as seis de
baixo, as seis restantes, que de facto partem os $18$ lados --- reaparecem, e
os $35$ mapas de retorno são todos rotações em torno de $O$, com
$\Gamma=C_3$. Os dois hexaflexágonos clássicos dão portanto a \emph{mesma}
resposta, e é o máximo que o Corolário~\ref{cor:hex} permite.
\end{teorema}

O Corolário~\ref{cor:hex} tem uma hipótese --- que o brinquedo não se abra
para fora do hexágono --- e ela merece ser testada, não assumida.
Implementámos por isso o análogo triangular do modelo com camadas
(\texttt{hexcamadas.py}): estados dobrados sobre o reticulado $\Z[\omega]$,
com o movimento geral ``refletir numa reta do reticulado um bloco contíguo de
camadas cujas charneiras cortadas assentem todas nessa reta'' --- que inclui
as dobras parciais, as aberturas, e em particular as aberturas \emph{por um
lado exterior}, isto é, as tentativas de fuga. A órbita do tri-hexaflexágono
tem $402$ estados; os seus contornos descem até um único triângulo (tamanhos
$1,2,3,6$ com $216,144,36,6$ estados), de modo que o modelo \emph{permite} a
fuga; e ainda assim todos os mapas de retorno alcançáveis são rotações em
torno de $O$. A hipótese do corolário não é uma restrição imposta ao
brinquedo: é uma consequência das camadas.

\begin{figure}[t]
\centering
\includegraphics[width=\linewidth]{fig/pinca.png}
\caption{A flexão de pinça, o único flex de um hexaflexágono. (a) pinçam-se
três raios alternados --- é uma dobra \emph{simultânea}, e é por isso que um
modelo de uma reta de cada vez não a alcança; (b) o brinquedo fica no
tri-dobrado, três ladrilhos alternados, cada um com duas folhas por ladrilho da
face; (c) abre-se pelos outros três raios, e cada braço manda um bloco
contíguo de camadas para o ladrilho vazio ao lado; (d) o hexágono volta, com as
faces trocadas. Todas as retas usadas passam pelo centro --- é o
Teorema~\ref{thm:retas} em ação.}\label{fig:pinca}
\end{figure}

Um modelo de uma reta de cada vez não realiza, porém, a \emph{flexão de
pinça}, que é uma dobra \emph{simultânea} por três raios alternados --- e o
tri-dobrado de três ladrilhos alternados não se obtém por dobras sucessivas.
Implementámo-la à parte (\texttt{hexcamadas.pinca}: os três pares, montanha ou
vale em cada um, e a abertura simultânea pelos outros três raios). Com ela a
órbita do tri-hexaflexágono passa de $402$ para $6756$ estados, dos quais
$696$ hexagonais --- e \emph{todos} os $35$ mapas de retorno continuam a ser
rotações em torno de $O$. É a confirmação do Teorema~\ref{thm:retas} com a
flexão verdadeira, e não apenas com o seu substituto.

Isto bate com o que Gardner registou em 1956 \cite{gardner}: cada face do
hexa-hexa aparece em mais de uma forma, ``devido a uma rotação dos triângulos
componentes uns em relação aos outros'', e as seis faces exibem ao todo $18$
configurações --- $6\times3$, que é exatamente $|\Gamma|=3$ arranjos por face.
É também a correção de uma leitura demasiado generosa: contar os mapas de
retorno de \emph{todas} as dobraduras combinatórias atribuía aos
hexaflexágonos um $\Gamma$ de posto $2$, quando o Teorema~\ref{thm:retas}
mostra que nenhum mapa de posto positivo é sequer realizável, porque toda
dobra de um hexágono de seis triângulos passa pelo centro.


\section{Os quatro flexágonos clássicos}\label{sec:quatro}

O mesmo método --- enumerar as dobraduras planas do plano do brinquedo sobre a
sua face, ler os \emph{mapas de retorno} e calcular $\Gamma$ --- resolve os
quatro flexágonos clássicos. Para os hexaflexágonos a tira é reta, com o
último triângulo colado no primeiro, e as \emph{faces não são postuladas}: uma
face é um conjunto de seis lados de célula que ocupa os seis ladrilhos, e os
estados que exibem o mesmo conjunto dão os mapas de retorno.

\begin{teorema}\label{thm:quatro}
Com $O$ o centro da face e $\Lambda_\square=2\Zi$:
\begin{center}
\renewcommand{\arraystretch}{1.3}
\begin{tabular}{lllll}
\toprule
flexágono & face & $\Gamma$ & soluções & inteiras?\\
\midrule
tri-tetraflexágono & quadrado & $2\Z$ (posto $1$) & $h(e^{i\pi z})$ & sim\\
hexa-tetraflexágono & quadrado & $\Lambda_\square$ (posto $2$, $|P|=1$) & $\C(\wp,\wp')$ & \textbf{não}\\
tri-hexaflexágono & hexágono & $C_3$ (posto $0$, $|P|=3$) & $h\big((z-O)^3\big)$ & sim\\
hexa-hexaflexágono & hexágono & $C_3$ (posto $0$, $|P|=3$) & $h\big((z-O)^3\big)$ & sim\\
\bottomrule
\end{tabular}
\end{center}
Só o hexa-tetraflexágono força polos: é o único dos quatro cujo $\Gamma$ tem
posto $2$, e a sua curva $\C/2\Zi$ é a lemniscática, $j=1728$.
\end{teorema}

\begin{figure}[t]
\centering
\includegraphics[width=\linewidth]{fig/quatro_faces.png}
\caption{Uma solução para cada flexágono, no retrato de fase da sua face. Nas
duas faces quadradas vê-se a diferença de posto: $\cos$ é periódica num só
eixo, $\wp$ é duplamente periódica e tem polos (os pontos onde todas as cores
se encontram). Nas duas faces hexagonais a simetria de ordem três em torno do
centro é imediata --- é literalmente $\Gamma=C_3$ a olho nu ---, e ambas são
soluções \emph{inteiras} ou quocientes delas.}\label{fig:quatrofaces}
\end{figure}

Os números: o anel de $12$ quadradinhos admite $70$ dobraduras planas sobre o
quadrado $2\times2$ e $216\,768$ estados dobrados; o octomino rasgado e colado
do tri-tetraflexágono, $6$ dobraduras; a tira de $10$ triângulos, $6$; a de
$19$ triângulos, $1362$. O que separa as quatro linhas da tabela não é o
tamanho do brinquedo mas a geometria da sua face, pelo
Teorema~\ref{thm:retas}: no quadrado há retas de dobra que não passam pelo
centro e o grupo pode ter posto $2$; no hexágono não há, e o grupo é
forçosamente finito. Dentro do caso quadrado, o que separa o tri do hexa é
quantos arranjos a flexão realmente alcança --- Teorema~\ref{thm:decidido}.


\section{A caracterização combinatória}\label{sec:paridade}

O Teorema~\ref{thm:decidido} foi obtido por busca exaustiva em $216\,768$
estados dobrados. Esta seção mostra que a busca é dispensável: o arranjo de
\emph{cada} face é função de quatro dígitos ternários.

\begin{teorema}[fórmula de paridade]\label{thm:paridade}
Numa dobradura plana do plano de um tetraflexágono sobre a face $2\times2$,
seja $\kappa_L\in\{0,1\}$ a coluna que a folha $L$ ocupa na face e $c_L$ a
coluna que ela ocupa no plano; analogamente $\rho_L$ e $r_L$ para as linhas.
Então a parte de rotação e o lado visível de cada folha são
\[
a_L\equiv 2(\kappa_L+c_L)\ (\mathrm{mod}\ 4),\qquad
s_L\equiv(\kappa_L+c_L)+(\rho_L+r_L)\ (\mathrm{mod}\ 2)\ \ (\text{a menos de constante}).
\]
\end{teorema}

\begin{proof}
As retas de dobra são $x\in\Z$ e $y\in\Z$, logo $W\cong D_\infty\times D_\infty$
age separadamente em $x$ e em $y$ e cada $g_L$ se escreve $(u_L,v_L)$. Uma
isometria $(u,v)$ é uma rotação de $180^\circ$ exatamente quando $u$ e $v$ são
ambas reflexões, e é direta quando ambas são do mesmo tipo; daí
$a_L=2\,t_x(L)$ e $s_L=t_x(L)+t_y(L)$, onde $t_x=1$ se $u_L$ inverte. Enfim,
se $u_L$ é translação a célula $[c,c+1]$ vai para $[c+2m,c+1+2m]$ e
$\kappa\equiv c$; se inverte, vai para $[-c-1+2m,-c+2m]$ e
$\kappa\equiv c+1$. Logo $t_x\equiv\kappa+c$.
\end{proof}

\begin{corolario}[fórmula dos intervalos]\label{cor:intervalos}
Para duas peças $j,j'$ de uma face exibida,
\[
a_j-a_{j'}\;\equiv\;2\cdot\#\{\text{charneiras \emph{verticais} dobradas entre
as duas peças ao longo da tira}\}\pmod 4 .
\]
O vetor de rotações de um arranjo lê-se dos bits de dobra: nenhuma isometria é
necessária.
\end{corolario}

\begin{proposicao}[uma charneira solta por bloco]\label{prop:solta}
As doze charneiras do anel repartem-se em quatro blocos de três paralelas,
$(0,1,2)$, $(3,4,5)$, $(6,7,8)$, $(9,10,11)$. Dobrar quatro células numa janela
de duas obriga a que \emph{exatamente uma} charneira de cada bloco fique
solta. Uma dobradura é portanto um ponto
$u=(u_1,u_2,u_3,u_4)\in\{0,1,2\}^4$, e das $81$ possibilidades ocorrem $25$.
Um flex vertical ou fixa $(u_1,u_3)$ ou muda \emph{as duas} coordenadas (as
charneiras soltas assentam no vinco, as dobradas na aresta livre, e a abertura
corta uma das duas famílias); um flex horizontal age do mesmo modo em
$(u_2,u_4)$.
\end{proposicao}

\begin{teorema}[o arranjo é função de $u$]\label{thm:combinatoria}
Seja $u$ uma das $25$ coordenadas. Para cada face $k$:
\begin{itemize}\itemsep2pt
\item a peça $j$ ocupa a posição $(\kappa_j,\rho_j)$, onde $\kappa_j$ é a
      paridade do número de charneiras verticais \emph{soltas} entre a peça $1$
      e a peça $j$ ao longo da tira, e $\rho_j$ a mesma coisa com as
      horizontais;
\item $k$ é exibível em $u$ se e só se essas quatro posições são distintas;
\item nesse caso a assinatura de rotações é a do
      Corolário~\ref{cor:intervalos}.
\end{itemize}
O conjunto de arranjos assim previsto coincide \emph{exatamente}, face a face,
com o que a enumeração dos $216\,768$ estados dobrados produz
(\texttt{paridade.py}; verificado em \texttt{verify\_camadas.py}).
\end{teorema}

Para o hexa-tetraflexágono a leitura é imediata:

\begin{center}
\renewcommand{\arraystretch}{1.2}
\begin{tabular}{lll}
\toprule
face & exibível em $(u_1,u_3)$ & assinatura\\
\midrule
$3,4,5,6$ & várias coordenadas & \emph{única}: $(0,0,0,0)$\\
$1$ & $(0,0),(1,1)$ & $(0,2,0,2)$\\
$1$ & $(2,2)$ & $(0,0,0,0)$\\
$2$ & $(1,1),(2,2)$ & $(0,2,0,2)$\\
$2$ & $(0,0)$ & $(0,0,0,0)$\\
\bottomrule
\end{tabular}
\end{center}

\begin{corolario}
$|P|=1$ --- isto é, $\Gamma$ é um reticulado puro --- se e só se nenhuma face é
exibida em duas coordenadas de assinaturas diferentes. As faces $3$ a $6$ têm
assinatura única e não dão trabalho. As faces $1$ e $2$ só se exibem na
\emph{diagonal} $u_1=u_3$, e aí cada uma tem duas assinaturas.
\end{corolario}

Resta portanto um único bit, e ele tem uma regra --- e a regra tem uma
simetria que vale a pena enunciar primeiro. As quatro charneiras soltas
(uma por bloco) cortam o anel em \emph{quatro arcos}, um por posição: todo
estado dobrado é um ciclo de quatro \emph{pats}, e os seus tamanhos são
$3+u_{i+1}-u_i$. Nas três coordenadas ``simétricas'' $u_1=u_2=u_3=u_4$ os
quatro pats têm três folhas cada, e a folha de topo de um pat pode ser a
primeira, a do meio ou a última do seu arco: $81$ padrões, e todos os $81$
existem como estado dobrado válido. Mas a flexão não os alcança todos.

\begin{figure}[t]
\centering
\includegraphics[width=.92\linewidth]{fig/extremos.png}
\caption{A regra dos extremos em $u=(2,2,2,2)$. Cada uma das quatro posições
guarda um pat de três folhas; a laranja é a \emph{folha de saída} do pat (a que
precede a charneira solta), e a de cima de cada pilha é a que se vê. À esquerda,
um dos oito padrões que a flexão alcança: a saída está no topo em dois pats
opostos. À direita, o que a face $1$ exigiria --- a saída no topo nos quatro ---
e que nunca ocorre. As setas dão a ordem em que a tira percorre as quatro
posições. (As camadas do meio estão desenhadas em ordem arbitrária: só o topo
importa.)}\label{fig:extremos}
\end{figure}

\begin{teorema}[regra dos extremos]\label{thm:extremos}
Nas coordenadas simétricas:
\begin{itemize}\itemsep2pt
\item $u=(1,1,1,1)$ --- charneiras soltas \emph{ao meio} de cada bloco ---
      a flexão alcança os $81$ padrões;
\item $u=(2,2,2,2)$ alcança exatamente $8$: a \emph{última} folha do arco está
      no topo em exatamente dois pats \emph{opostos};
\item $u=(0,0,0,0)$ alcança exatamente $8$: o mesmo com a \emph{primeira}
      folha --- o enunciado espelhado, como tem de ser, porque inverter o
      sentido do anel troca $(0,0,0,0)$ com $(2,2,2,2)$.
\end{itemize}
\end{teorema}

É isto que fecha o problema. A assinatura excepcional da face $1$ vive só em
$u=(2,2,2,2)$, e exibi-la aí exigiria os topos $(2,5,8,11)$, isto é, a última
folha nos \emph{quatro} pats --- e a regra permite no máximo dois. A face $6$
pede $(2,3,8,9)$, última--primeira--última--primeira: dois opostos, permitido;
a face $4$ pede $(0,5,6,11)$, o padrão complementar; ambas ocorrem. E, no
espelho, a assinatura excepcional da face $2$ vive em $u=(0,0,0,0)$ e é morta
pelo mesmo argumento. Cada face fica assim com uma \emph{única} assinatura ao
longo de toda a órbita --- a face $1$ aparece em $u=(0,1,0,1)$, $(1,0,1,0)$ e
$(1,1,1,1)$, a face $2$ em $(1,1,1,1)$, $(1,2,1,2)$ e $(2,1,2,1)$ --- donde
$\varepsilon$ constante, $|P|=1$ e $\Gamma=2\Zi$, sem nenhuma busca.

\subsection*{Por que a regra vale: o grafo de flexão é desconexo}

A regra dos extremos não é um acidente da órbita: é a sombra de um facto
estrutural. O grafo de flexão sobre os $216\,768$ estados dobrados
\emph{não é conexo}. O flexágono montado vive numa componente de $19\,200$
estados, que exibe as seis faces. Já um estado dobrado que exiba o arranjo
excepcional da face $1$ --- e ele existe, é um estado dobrado perfeitamente
válido --- vive numa componente de \emph{vinte} estados, disjunta daquela, em
que nenhuma outra face chega a aparecer: um objeto de papel rígido, preso em
$u=(2,2,2,2)$, que exibe a face $1$ e mais nada. Para passar dele ao flexágono
seria preciso desdobrar o anel por completo e voltar a dobrá-lo; nenhuma
sequência de flexões o faz.

E a paisagem toda é assim --- calculámo-la por inteiro.

\begin{figure}[t]
\centering
\includegraphics[width=.86\linewidth]{fig/grafo_tri.png}
\caption{O grafo de flexão do tri-tetraflexágono \emph{inteiro} --- ele é
pequeno o bastante para caber na página. São $24$ classes de estados dobrados
(a menos das simetrias globais da face), $12$ flexões e $14$ componentes. A
primeira é o brinquedo: um triângulo, e nos seus três vértices estão as três
faces $F$, $V$, $E$. As outras treze são grãos, e uma delas exibe a face $V$
num arranjo que nenhuma flexão alcança --- o mesmo fenômeno que a
Figura~\ref{fig:grafo} mostra em escala grande no
hexa-tetraflexágono.}\label{fig:grafotri}
\end{figure}

\begin{figure}[t]
\centering
\includegraphics[width=\linewidth]{fig/grafo_componentes.png}
\caption{O grafo de flexão. Nos três primeiros painéis, componentes inteiras
--- grãos de $4$, $8$ e $20$ estados, desenhados como grafos: são objetos de
papel que dobram mas não flexionam. No último, as $20\,405$ componentes em
escala logarítmica dupla: um pó cujo maior grão tem $200$ estados, e, isolada à
direita, a ilha do flexágono com $19\,200$.}\label{fig:grafo}
\end{figure}

\begin{teorema}[classificação das componentes]\label{thm:componentes}
O grafo de flexão dos $216\,768$ estados dobrados do anel de doze
quadradinhos sobre o quadrado $2\times2$ tem exatamente $20\,405$
componentes: \emph{uma} com $19\,200$ estados --- o flexágono --- e
$20\,404$ com $200$ estados ou menos. A lista completa dos tamanhos é
\begin{center}\small
\renewcommand{\arraystretch}{1.1}
\begin{tabular}{r|rrrrrrrrrr}
tamanho & $4$ & $8$ & $12$ & $16$ & $20$ & $24$ & $28$ & $32$ & $36$ & $40$\\
quantas & $9244$ & $5288$ & $2208$ & $1256$ & $640$ & $820$ & $402$ & $84$ & $48$ & $96$\\[2pt]
tamanho & $44$ & $48$ & $52$ & $56$ & $64$ & $72$ & $100$ & $128$ & $140$ & $200$\\
quantas & $32$ & $120$ & $8$ & $48$ & $32$ & $48$ & $8$ & $16$ & $2$ & $4$\\
\end{tabular}
\end{center}
(mais a ilha, $19\,200\times1$; a soma dá $216\,768$). Todos os tamanhos são
múltiplos de quatro, porque as quatro simetrias globais que preservam a órbita
agem livremente.
\end{teorema}

\begin{corolario}[critério de decisão]\label{cor:criterio}
Um estado dobrado é um estado do flexágono se e só se a sua componente tem
mais de $200$ estados. O teste termina em no máximo $201$ expansões: cresça-se
a componente e pare-se ao atingir $201$.
\end{corolario}

O espaço dos estados dobrados do anel é, portanto, \emph{uma ilha grande num pó
de grãos rígidos}: o flexágono é a ilha --- noventa e seis vezes maior do que o
maior dos grãos ---, e cada grão é um objeto de papel que se dobra mas não
flexiona.

É esta a razão de fundo. O quinto arranjo da face $1$ não é ``difícil de
alcançar'': está noutro mundo. A cadeia completa fica: $25$ coordenadas em vez
de $216\,768$ estados; a exibibilidade de cada face por paridades; a sua
assinatura pela fórmula dos intervalos; e, nas duas coordenadas extremas onde
as duas assinaturas colidiriam, a regra dos extremos, que é o critério de
componente. A enumeração exaustiva serve agora só de verificação
(\texttt{verify\_camadas.py} e \texttt{verify\_geral.py}, $65$ testes exatos).


\section{Os reticulados que ocorrem}\label{sec:reticulados}

A classificação do Teorema~\ref{thm:classificacao} deixa em aberto \emph{quais}
reticulados $\Lambda\subset\frac12\Zi$ efetivamente ocorrem. A resposta é
surpreendentemente rígida.

\begin{teorema}\label{thm:reticulados}
Entre os $5656$ flexes de quadrantes com $\Gamma_\Phi$ de posto $2$ ocorrem
exatamente $86$ reticulados distintos. Eles se distribuem em apenas
\emph{quatro} classes de homotetia, a saber $\tau\in\{i,\,2i,\,\frac{-1+2i}2,\,
\frac{-1+3i}2\}$, com
\[
j=1728,\qquad j=287496,\qquad j\approx-89{,}5950,\qquad j\approx-11663{,}3967,
\]
correspondentes às ordens de condutor $1,2,4,6$ em $\Zi$ (discriminantes
$-4,-16,-64,-144$). Além disso:
\begin{enumerate}
\item[(a)] se $|P_\Phi|=4$ --- o que ocorre em $5240$ dos $5656$ casos --- então
$\Lambda$ é um ideal gaussiano, $\Lambda=\alpha\Zi$, e portanto
$\C/\Lambda\cong\C/\Zi$ \emph{sempre}: a curva é a lemniscática, $j=1728$. As
normas que ocorrem são
$N(\alpha)\in\{1,2,4,5,8,10,16,18\}$ (em unidades de $\covol\Zi$);
\item[(b)] as outras três classes de homotetia só aparecem quando
$|P_\Phi|\le2$, em $8$ reticulados;
\item[(c)] em todos os casos $\C/\Lambda$ tem multiplicação complexa por uma
ordem de $\Q(i)$, como já garantia o Corolário~\ref{cor:cm}.
\end{enumerate}
\end{teorema}

\begin{proof}
A parte numérica é a enumeração exata (Seção~\ref{sec:computacional}); só o
item (a) precisa de argumento. Se $|P_\Phi|=4$ então $i\Lambda=\Lambda$, isto
é, $\Lambda$ é um $\Zi$-submódulo de $\Q(i)$ finitamente gerado, ou seja um
ideal fracionário; como $\Zi$ é principal, $\Lambda=\alpha\Zi$. A homotetia
$z\mapsto z/\alpha$ leva $\Lambda$ em $\Zi$.
\end{proof}

\begin{observacao}
Em particular a curva lemniscática não é apenas o caso ``genérico'': ela é o
\emph{único} caso possível assim que o flex gira alguma peça de $90^\circ$.
Os dois flexágonos de \cite{manual} caem no caso $\alpha=2$, $|P|=1$.
\end{observacao}

\section{Catálogos}\label{sec:catalogo}

O programa \texttt{catalogo.py} materializa a Seção~\ref{sec:gerador}: para cada
flexágono ele produz um diretório
\begin{center}
\texttt{funcoes-flexionaveis-<nome>/}\quad com\quad
\texttt{001/face1.png}\,\ldots\,\texttt{faceN.png},\ \texttt{mecanica.png},\
\texttt{formula.tex}
\end{center}
e um \texttt{catalogo.pdf} de duas colunas em que cada item traz o seu número,
as fórmulas das faces e o \emph{diagrama de flexão com os retratos de fase
encaixados}: no lugar dos rótulos $1a,1b,1c,1d$ do diagrama original aparecem os
quadrantes da figura, cada caixa no arranjo em que aquela face de fato se
apresenta naquele estado. As caixas cheias e tracejadas, as setas e o símbolo
$\circlearrowright$ reproduzem o diagrama de \cite{enimar}. As imagens
\texttt{faceN.png} são exatamente o que o gerador de planos de \cite{manual}
espera como entrada, de modo que cada item do catálogo é um flexágono
imprimível.

Cada função é sorteada na parametrização completa do
Teorema~\ref{thm:classificacao} e \emph{testada} antes de ser desenhada:
o resíduo $\max_j\|f\circ A_1^{-1}-f\circ A_j^{-1}\|$ é medido contra todos os
mapas de retorno e o item só é aceito abaixo de $10^{-7}$ (na prática obtém-se
$10^{-14}$). Assim o catálogo não é uma coleção de exemplos: é uma amostragem
certificada da família completa.

\section{Faces com mais de um protótipo}\label{sec:multiproto}

Se a face é ladrilhada por formas diferentes, um flex tem de respeitar as
classes de congruência: $\sigma$ permuta ladrilhos dentro de cada classe e
$\varepsilon_j$ percorre a co-classe das rotações que levam $T_j$ em
$T_{\sigma(j)}$. O grupo dos flexes deixa de ser um produto entrelaçado ---
é um subgrupo do produto directo dos produtos entrelaçados de cada classe ---
mas o Teorema~\ref{thm:principal} não muda: a sua demonstração só usa o
princípio da identidade ao longo das arestas, e não a congruência dos
ladrilhos. O que muda é a aritmética.

\begin{teorema}\label{thm:multiproto}
Seja $D$ uma face ladrilhada por protótipos com grupos de rotação
$C_{n_1},\dots,C_{n_r}$ e seja $\Phi$ um flex. Então exatamente um dos
seguintes casos ocorre.
\begin{enumerate}\itemsep2pt
\item Os $A_j$ têm um ponto fixo comum $p$. Então $\Gamma_\Phi=C_N$ é
      \emph{finito}, $N$ a ordem do grupo de pontos, e as soluções são
      $f(z)=h\big((z-p)^N\big)$, $h$ meromorfa --- entre elas funções
      inteiras não constantes. Aqui $N$ é arbitrário: a restrição
      cristalográfica \emph{não} se aplica, porque não há translações.
\item $\Gamma_\Phi$ contém uma translação não nula e $N\in\{1,2,3,4,6\}$:
      $\Gamma_\Phi$ é discreto e cai numa das nove classes do
      Teorema~\ref{thm:universal}.
\item $\Gamma_\Phi$ contém uma translação não nula e $N\notin\{1,2,3,4,6\}$:
      então $\Gamma_\Phi$ é \emph{denso} em $\C\rtimes U(1)$ e as únicas
      soluções são as constantes.
\end{enumerate}
\end{teorema}

\begin{proof}
$\Gamma_\Phi\le\C\rtimes U(1)$ e o seu grupo de pontos $P$ é gerado pelos
$\varepsilon_k^{-1}\varepsilon_j$, logo é cíclico finito de ordem $N$ dividindo
$\mathrm{lcm}(n_1,\dots,n_r)$. Se todos os $A_j$ fixam $p$, então
$\Gamma_\Phi$ fixa $p$ e é um grupo de rotações em torno de $p$: é $C_N$, e a
função invariante genérica é $h((z-p)^N)$ (Hauptmodul da órbifold
$\C/C_N$). Caso contrário $\Gamma_\Phi$ contém a translação não nula
$T=A_k^{-1}A_j$ para um par sem ponto fixo comum, e o subgrupo de translações
$\Lambda$ é invariante por $P$, isto é, é um $\Z[\zeta_N]$-módulo não nulo.
Se $N\in\{1,2,3,4,6\}$, $\Z[\zeta_N]$ é discreto e recai-se no
Teorema~\ref{thm:universal}. Se $N\notin\{1,2,3,4,6\}$, $\Z[\zeta_N]$ é denso
em $\C$, logo $\Gamma_\Phi$ é denso; uma função meromorfa invariante por um
grupo denso é constante (a órbita de um ponto regular acumula, e o princípio
da identidade fecha o argumento).
\end{proof}

\begin{figure}[t]
\centering
\includegraphics[width=.55\linewidth]{fig/roda.png}
\caption{A \emph{roda}: dois quadrados e três triângulos equiláteros fecham os
$360^\circ$ em torno de um vértice. A rotação de $150^\circ$ em torno desse
vértice troca os dois quadrados e fixa cada triângulo; como
$\langle150^\circ\rangle=C_{12}$ e todos os $A_j$ fixam o vértice, o grupo é
finito de ordem $12$ --- fora da lista cristalográfica, o que só é possível
porque não há translações. As soluções são $h\big((z-p)^{12}\big)$, e há
inteiras.}\label{fig:roda}
\end{figure}

\begin{exemplo}
Um quadrado unitário ($C_4$) com um triângulo equilátero colado em cima e
outro em baixo ($C_3$): a face tem três ladrilhos e dois protótipos.
Aritmética exata em $\Z[\zeta_{12}]$ (\texttt{multiproto.py}) dá $72$ flexes:
$32$ deles têm $\Gamma$ \emph{denso} --- só constantes ---, $29$ são
cristalográficos e $11$ são cíclicos finitos $C_N$ com
$N\in\{1,2,3,4,6\}$. Misturar protótipos torna a face \emph{mais} rígida, não
menos.
\end{exemplo}

\begin{corolario}
Num flexágono de verdade o caso denso nunca ocorre: pelo
Teorema~\ref{thm:retas}, $\Gamma\le W^+(D)\le\mathrm{Sym}^+(\mathcal T)$, que é
um grupo cristalográfico. A rigidez total é um fenómeno dos flexes
\emph{abstratos}.
\end{corolario}



\begin{figure}[t]
\centering
\includegraphics[width=\linewidth]{fig/variedade.png}
\caption{Oito itens do catálogo dos hexaflexágonos. Todas são
$f(z)=h(z^3)$ com $h$ racional sorteada, e todas exibem a olho nu a simetria de
ordem três em torno do centro que o Corolário~\ref{cor:hex} obriga. Nenhuma se
quebra, seja qual for a flexão de pinça.}\label{fig:variedade}
\end{figure}

Para os hexaflexágonos o catálogo é outro e muito mais simples, porque
$\Gamma=C_3$: as soluções são $f(z)=h\big((z-O)^3\big)$ com $h$ meromorfa
qualquer, e há inteiras. \texttt{catalogo\_hex.py} gera-as, certifica
$f(\omega z)=f(z)$ com resíduo $\sim10^{-14}$, desenha o retrato de fase do
hexágono, recorta os seis triângulos e monta o
\texttt{catalogo.pdf} de duas colunas, um para cada hexaflexágono. Os quatro
catálogos existem e compilam: $12$ itens para o tri-tetraflexágono, $8$ para o
hexa-tetraflexágono, $12$ para cada hexaflexágono. Com a semente $2026$ eles se
reproduzem exatamente.

\section{Observações finais}

\textbf{1. Uma leitura em termos de órbifolds.} O Teorema~\ref{thm:principal}
diz que uma função flexionável é uma função na órbifold $\C/\Gamma_\Phi$. As
oito linhas da tabela do Teorema~\ref{thm:classificacao} são exatamente as
órbifolds quocientes possíveis: $\C$, $\C(n)$, $\C^\ast$, o
disco-com-dois-pontos-cônicos, o toro $\C/\Lambda$, a ``almofada''
$S^2(2,2,2,2)$ e $S^2(2,4,4)$. Os dois flexágonos de \cite{manual} caem nos dois
casos mais simples: $\C^\ast$ para o tri (donde $h(e^{i\pi z})$) e o toro
$\C/2\Zi$ para o hexa (donde as elípticas). A Seção~\ref{sec:gerador}
transforma essa observação num algoritmo.

\smallskip
\textbf{2. Por que $\Zi$?} A resposta aritmética --- multiplicação complexa por
$\Q(i)$ --- é forçada por dois fatos elementares: os quadradinhos são
\emph{quadrados} (donde $\varepsilon_j\in\mu_4$) e os seus centros estão em
$\frac12\Zi$. Num ladrilhamento triangular o mesmo argumento produziria
$\mu_6$ e $\Z[\zeta_3]$ --- a curva \emph{equianarmônica} $g_2=0$, $j=0$ ---, e
é isso que se vê no censo dos $524\,880$ flexes triangulares abstratos; mas
nenhum hexaflexágono real lá chega, porque o Teorema~\ref{thm:retas} o confina
às rotações em torno do centro.
Seria interessante desenvolver esse caso.

\smallskip
\textbf{3. Uma sugestão prática.} Para montar um flexágono ``analítico'' basta
alimentar o gerador de \cite{manual} com as imagens produzidas por
\texttt{catalogo.py}: seis funções elípticas distintas de reticulado $2\Zi$ para
o hexa, três funções de $e^{i\pi z}$ para o tri. Cada face é um desenho
diferente e nenhuma se quebra, seja qual for a flexão.

\smallskip
\textbf{4. As três perguntas que ficaram em aberto.}
\emph{(a) Quais dobraduras são realizáveis, e quais são alcançáveis?}
Resolvido. A realizabilidade das camadas (condições de Justin) não separa nada:
as $70$ dobraduras planas do hexa-tetraflexágono passam-na todas. A
alcançabilidade, sim: no modelo com camadas do §\ref{sec:camadas} a órbita do
estado montado contém $19\,200$ dos $216\,768$ estados dobrados, e nela as
faces $1$ e $2$ exibem-se em três arranjos e as faces $3$ a $6$ em dois ---
exatamente o diagrama de \cite{enimar}. Os dois arranjos combinatórios a mais
são os únicos com $\varepsilon$ não constante, e são inalcançáveis: exigiriam
separar camadas não contíguas. O enquadramento $\C(\wp)\subseteq{}$soluções
$\subseteq\C(\wp,\wp')$ colapsa assim no extremo de cima:
$\Gamma=2\Zi$ é um reticulado puro, $|P|=1$, e as soluções são
\emph{todas} as funções elípticas de $2\Zi$ (Teorema~\ref{thm:decidido}).
Como a órbita foi semeada com todas as ordens de camadas do estado montado, a
exclusão é incondicional.

\smallskip
\emph{(b) O grupo de flexão.} Nos dois tetraflexágonos todos os mapas de
retorno têm $\varepsilon$ constante --- são translações a menos de uma
isometria global ---, e nos dois hexaflexágonos são rotações em torno do
centro. Isso não é um acidente dos dois brinquedos: é o
Teorema~\ref{thm:retas}, que dá $\Gamma\le W^+(D)$ para qualquer flexágono de
face $D$. $\Gamma$ é invariante por pós-composição com isometrias globais
(Lema~\ref{lem:rho}), pelo que não distingue estados que difiram por uma
rotação do brinquedo --- é exatamente o quociente certo para a pergunta
analítica.

\smallskip
\emph{(c) O caso triangular.} Resolvido, e com uma inversão: o
\emph{ladrilhamento} triangular dá $\mu_6$, $\Z[\zeta_3]$ e a curva
equianarmônica --- e o censo dos $524\,880$ flexes abstratos mostra que
$99{,}69\%$ deles forçam polos ---, mas os hexaflexágonos \emph{reais} não
usam nenhum desses flexes. Como toda reta de dobra de um hexágono de seis
triângulos passa pelo centro, o seu $\Gamma$ é finito, e igual a $C_3$ nos
dois hexaflexágonos clássicos, com soluções
$h\big((z-O)^3\big)$ (Teorema~\ref{thm:hexaflex}), em acordo com as $18$
configurações que Gardner contou. O contraste com o caso quadrado continua a
ter uma causa única, mas é outra que não o ladrilho: é o facto de o quadrado
poder ser aberto por um lado e o hexágono não.

\smallskip
\textbf{5. E o que ficou de fora.} As duas perguntas que a versão anterior
deixou pendentes estão fechadas. As faces com mais de um protótipo são o
Teorema~\ref{thm:multiproto}: o teorema de caracterização não muda, a
restrição cristalográfica passa a valer só quando há translações, e abre-se um
caso novo --- $\Gamma$ denso, só constantes --- que nenhum flexágono realiza.
E o cálculo de $W^+(D)$ deixou de ser feito à mão: \texttt{wgrupo.py} fá-lo
para qualquer polióminos ou poliamante (Observação~\ref{obs:wgrupo}), de modo
que uma face nova se resolve em três passos --- calcular $W^+(D)$, que fixa as
classes possíveis; correr o modelo com camadas, que escolhe dentro delas; ler
o Hauptmodul na tabela do Teorema~\ref{thm:universal}.

A caracterização combinatória pedida está na Seção~\ref{sec:paridade}, e está
completa: o arranjo de cada face é função dos quatro dígitos $u$
(Teorema~\ref{thm:combinatoria}), e o único bit que os dígitos não fixam é
decidido pela regra dos extremos (Teorema~\ref{thm:extremos}). Nada disto
usa a busca exaustiva --- ela serve agora só de verificação.

A classificação das componentes está feita (Teorema~\ref{thm:componentes}) e
com ela o critério de decisão do Corolário~\ref{cor:criterio}; a regra dos
extremos é agora um corolário verificado dela. Vale a pena dizer com precisão
o que isso deixa, e o que não deixa, em aberto.

Em aberto não fica a \emph{classificação}: a lista das $20\,405$ componentes é
completa, e pertencer à ilha decide-se em no máximo $201$ expansões. O que
falta é uma \emph{fórmula} onde hoje temos uma \emph{busca} --- uma questão de
forma, e não de substância. Fica registado, porém, que ela não virá do lugar
onde seria natural procurá-la.

\begin{proposicao}\label{prop:naotopologico}
Nenhum invariante topológico do mergulho distingue duas componentes: todos os
estados dobrados do anel são ambientalmente isotópicos entre si.
\end{proposicao}

\begin{proof}
Todo estado dobrado plano de um polígono de papel é atingível a partir do papel
plano por um movimento contínuo sem auto-intersecção \cite{cfp}. Logo cada
estado dobrado do anel é isotópico ao anel plano e, por transitividade, dois
quaisquer são isotópicos entre si. Um invariante de isotopia é, portanto,
constante em todo o espaço dos estados, e não pode separar componentes.
\end{proof}

O que distingue a ilha dos grãos não é, pois, a \emph{forma} do objeto
dobrado --- é o \emph{movimento} permitido: desdobrar o anel por completo
levaria qualquer grão ao flexágono, e é exatamente isso que um flex não faz.
Um invariante teria de ler a restrição, não a geometria.

\begin{observacao}
Também não parece vir de uma estatística local. Testámos, na ilha e no pó, as
coordenadas $u$ (as $25$ ocorrem na ilha), o enrolamento da banda pelas camadas
$\sum\delta_i$ (todos os valores de $-4$ a $4$ ocorrem em ambos), o grau no
grafo de flexão (ilha $2$--$22$, pó $2$--$16$: sobrepõem-se), o perfil de
camadas por posição (os cinco perfis ocorrem em ambos), o sinal da permutação
de leitura e a paridade dos lados: nenhum separa. Isto é evidência, não
demonstração --- ao contrário da Proposição~\ref{prop:naotopologico}.
\end{observacao}

\begin{thebibliography}{9}
\bibitem{manual}
A.~Milanés Barrientos e G.~Ribeiro,
\emph{Papel, quadrados misteriosos e \ldots matemática}, Coleção
``Matemática fora da caixa'', Projeto Visitas no Mundo da Matemática,
Departamento de Matemática, UFMG, Belo Horizonte, 2025.
ISBN 978-65-01-47006-1.

\bibitem{enimar}
Projeto Visitas no Mundo da Matemática, \emph{Diagrama de flexão do
hexa-tetraflexágono}, comunicação do EnIMaR, UFMG, 2025.

\bibitem{demaine}
E.~D.~Demaine e J.~O'Rourke, \emph{Geometric Folding Algorithms: Linkages,
Origami, Polyhedra}, Cambridge University Press, 2007.

\bibitem{cfp}
E.~D.~Demaine, S.~L.~Devadoss, J.~S.~B.~Mitchell e J.~O'Rourke,
\emph{Continuous foldability of polygonal paper}, Proc.\ 16th Canadian
Conference on Computational Geometry (CCCG 2004), 64--67.

\bibitem{gardner}
M.~Gardner, \emph{Hexaflexagons and Other Mathematical Diversions},
University of Chicago Press, 1988.

\bibitem{pook}
L.~P.~Pook, \emph{Serious Fun with Flexagons: A Compendium and Guide},
Springer, 2009.

\bibitem{ahlfors}
L.~V.~Ahlfors, \emph{Complex Analysis}, 3ª ed., McGraw--Hill, 1979.

\bibitem{silverman}
J.~H.~Silverman, \emph{Advanced Topics in the Arithmetic of Elliptic Curves},
Springer GTM 151, 1994.

\bibitem{wegert}
E.~Wegert, \emph{Visual Complex Functions: An Introduction with Phase Portraits},
Birkhäuser, 2012.
\end{thebibliography}

\end{document}
